% Standalone source, assembled from the supplied modular sections.
% Regenerate after editing with: python3 build.py --assemble-only
\documentclass[11pt,a4paper]{article}
\usepackage[T1]{fontenc}
\usepackage{lmodern}
\usepackage[margin=27mm,headheight=14pt]{geometry}
\usepackage{amsmath,amssymb,amsthm,mathtools}
\usepackage{microtype,booktabs,longtable,array,xcolor,xurl}
\usepackage{enumitem,fancyhdr}
\usepackage[colorlinks=true,linkcolor=blue!45!black,citecolor=blue!40!black,urlcolor=blue!40!black]{hyperref}
\usepackage{bookmark}
\pagestyle{fancy}\fancyhf{}
\fancyhead[L]{\small Exact resonance and coordinate transport}
\fancyhead[R]{\small ProveIt research continuation}
\fancyfoot[C]{\thepage}
\setlength{\parskip}{3pt}
\setlength{\parindent}{1.25em}
\setlist{itemsep=3pt,topsep=5pt}
\setcounter{tocdepth}{2}
\emergencystretch=2em
\numberwithin{equation}{section}
\newtheorem{theorem}{Theorem}[section]
\newtheorem{proposition}[theorem]{Proposition}
\newtheorem{lemma}[theorem]{Lemma}
\newtheorem{corollary}[theorem]{Corollary}
\theoremstyle{definition}
\newtheorem{definition}[theorem]{Definition}
\theoremstyle{remark}
\newtheorem{remark}[theorem]{Remark}
\newtheorem{example}[theorem]{Example}
\DeclareMathOperator{\Li}{Li}
\DeclareMathOperator{\FP}{FP}
\DeclareMathOperator{\CT}{CT}
\DeclareMathOperator{\Res}{Res}
\DeclareMathOperator{\Log}{Log}
\DeclareMathOperator{\sgn}{sgn}
\DeclareMathOperator{\lcm}{lcm}
\newcommand{\dd}{\,\mathrm d}
\newcommand{\eps}{\varepsilon}
\newcommand{\C}{\mathbb C}
\newcommand{\R}{\mathbb R}
\newcommand{\Q}{\mathbb Q}
\newcommand{\Z}{\mathbb Z}
\newcommand{\N}{\mathbb N}
\newcommand{\src}[1]{\nolinkurl{#1}}
\newcommand{\pin}{\texttt{445f754610e1377939235794de84ed9075fc09f5}}
\title{\textbf{Exact Resonance and Coordinate Transport}\\[6pt]
\large Dougall--Hurwitz identities, nonlinear Stieltjes finite parts,\\
and polylogarithmic correlations at unequal frequencies}
\author{ProveIt research continuation\\
\small Prepared for Vladimir Reshetnikov with OpenAI ChatGPT}
\date{10 October 2026}
\hypersetup{pdftitle={Exact Resonance and Coordinate Transport},pdfauthor={Prepared for Vladimir Reshetnikov with OpenAI ChatGPT},pdfsubject={Dougall resonance, Stieltjes finite parts, harmonic identities, and colored polylogarithms}}

\begin{document}

% BEGIN sections/00_scope.tex
\maketitle
\begin{abstract}
We answer three explicit continuation questions in the current ProveIt
polylogarithm programme. A centered form of Dougall's very-well-poised
summation has only even inverse powers in its asymptotic expansion. Its
Hurwitz-subtracted endpoint admits a closed formula at every nonpositive
integer resonance, and every spectral derivative has a finite expression
in Gamma derivatives, Hurwitz-zeta derivatives, and generalized Stieltjes
constants. Reciprocal-Gamma zeros are handled before coefficient extraction.
A quartic central-binomial specialization gives an all-resonance family
with explicit harmonic constants; a weighted primitive reduces the
hypergeometric order and transports the subtraction terms to colored
polylogarithms.

A second theorem gives the exact change of a one-sided Hadamard finite
part under any analytic orientation-preserving local coordinate change.
The correction is a finite residue, obeys a composition law, and specializes
to all argument derivatives of all generalized Stieltjes functions.
For a coordinate change tangent to the identity, the correction for
$\gamma_n^{(r)}$ vanishes for $n\ge 2r$, and this threshold is sharp.
Finally, a finite multiplication formula reduces three unequal integer
frequencies to ordinary colored Tornheim kernels, with explicit lower
Stieltjes terms accounting for the fixed endpoint coordinate. It also
produces convergent log-Gamma correlations and negative-order reduction
identities. Every new assertion is supplied with a proof and accompanied
by reproducible finite algebra or independent numerical diagnostics.
\end{abstract}

\noindent\textbf{Status.} The contributions are relative to the inspected
repository and incoming reports. Classical summations and product-integral
methods are credited explicitly. Historical priority for every derived
special case is not asserted. The current Gaussian $S_6$ and $S_8$
identities remain conjectural.
\clearpage
\tableofcontents
\clearpage

\section{Research scope, source baseline, and conventions}
\label{sec:scope}

\subsection{The questions addressed}
The repository's consolidated manuscript and its incoming research reports
already contain an extensive calculus of polylogarithms, harmonic sums,
Hurwitz-zeta jets, and regularized Stieltjes products. A useful continuation
must retain that accumulated proof status. The baseline used here is the
commit
\begin{center}\small\pin.\end{center}
The two requested source locations are the manuscript under
\src{Analysis/Polylogarithms/docs/manuscript} and the archives under
\src{docs/incoming} \cite{repo}.

The present work concentrates on three questions explicitly posed in the
newest incoming material:
\begin{enumerate}
\item The Gauss--Hurwitz report asks for a fixed Dougall summation with a
complete resonant coefficient law and finite Stieltjes closure
\cite[Section 9, Question 1]{incoming-gh}.
\item The Mellin--dilation and Stieltjes--harmonic reports ask for the exact
nonlinear coordinate correction to finite parts, including compatibility
with composition \cite{incoming-md,incoming-sh}.
\item The Stieltjes--harmonic report asks for three unequal integer
frequencies and a finite colored-kernel description with the local
normalization included \cite[Further research: three unequal integer
frequencies]{incoming-sh}.
\end{enumerate}
The answer to each is constructive. The first uses even Bernoulli
coefficients and a Gamma-residue identity. The second uses a local residue
that depends on a finite coordinate jet. The third uses a finite Hurwitz
multiplication formula before taking any finite part. In the third case,
integer weights do not require a new transcendental kernel: finitely many
ordinary colored Tornheim kernels suffice.

\subsection{Reading and audit boundaries}
The manuscript directory, its editorial ledger, relevant depth,
integration, differentiation, and discovery sections, and all five incoming
archives were retrieved. The complete text of the four reports most
directly connected to the new identities was examined, and the fifth
report's scope and relevant source-status discussion were checked
\cite{incoming-nhj,incoming-continuation}.
The source manifest records paths and hashes; the five archive bytes were
checked against their Git blob identifiers. This is a targeted proof
audit, not a certification of every theorem in the 417-page source book.

The inherited results used below are rederived when the proof is short:
the six-sign-region Fourier formula, the elementary Stieltjes
multiplication law, and the local residue mechanism. The longer
independent proof of Dougall's summation is classical and cited
\cite[Eq.~16.4.9]{dlmf-hyper}. The existing affine correction is recovered
as a specialization, which checks the normalization without presenting
that earlier theorem as a fresh discovery.

\subsection{Classical antecedents}
Gamma-ratio asymptotics and generalized hypergeometric endpoint summation
are classical. B\"uhring's study of partial sums includes the zero-excess
very-well-poised constant and a quartic-binomial specialization
\cite[Section 4]{buhring}. In particular, the first quartic-binomial row
below is a recovered classical identity. The extension proved here is
the centered coefficient structure, the uniform nonpositive-integer
resonance law, and its complete spectral coefficient calculus with a
specified Hurwitz subtraction.

Hurwitz product integrals and Tornheim sums were studied systematically
by Espinosa and Moll \cite{em-hurwitz,em-tornheim}; colored Tornheim
relations have a separate literature \cite{zhao}. The Fourier and
multiplication formulas for Hurwitz zeta and the Gamma connection are
standard \cite{dlmf-hurwitz}. Our unequal-frequency result is a finite
transport theorem, with the fixed-coordinate terms written out. Its
finite expression does not imply arithmetic independence or a minimal
number of period coordinates.

Analytic regularization and Hadamard finite parts also have classical
antecedents. The local theorem below is proved directly so that its
Jacobian convention, residue sign, and composition rule can all be
checked without an unspecified regularization prescription. The
specialized Stieltjes coefficient formula and sharp high-index
invariance are the concrete additions needed by the source programme.

\subsection{Conventions}
All powers of positive real numbers use the real logarithm. The spectral
and argument derivatives are distinguished by
\begin{align}
 \zeta(1+u,a)&=\frac1u+\sum_{n\ge0}\frac{(-1)^n}{n!}\gamma_n(a)u^n,
 &\gamma_0(a)&=-\psi(a), \label{eq:laurent}\\
 \zeta^{(k)}(s,a)&=\frac{\partial^k}{\partial s^k}\zeta(s,a),
 &\gamma_n^{(r)}(a)&=\frac{\mathrm d^r}{\mathrm da^r}\gamma_n(a).
\end{align}
The same superscript on $\psi$ denotes an argument derivative. Generalized
harmonic numbers are $H_N^{(j)}=\sum_{k=1}^Nk^{-j}$, with $H_0^{(j)}=0$;
$H_N=H_N^{(1)}$. The rising factorial is $(x)_N$, and $[t^k]F$ denotes a
Taylor coefficient, not a derivative. Bell polynomials are normalized by
\[
 \exp\!\left(\sum_{j\ge1}x_j\frac{t^j}{j!}\right)
 =\sum_{k\ge0}Y_k(x_1,\ldots,x_k)\frac{t^k}{k!}.
\]
The symbol $e(x)$ means $\exp(2\pi i x)$ when introduced in the Fourier
sections. On the circle, every finite part uses the right local
coordinate at each singular point. The notation $\FP_x$ is reserved for
that prescribed constant term; it does not denote an unnormalized
analytic continuation.

The proofs justify limiting operations analytically. Exact symbolic
checks verify finite identities and their implementations. Numerical
checks compare independent representations and detect sign, scale, and
transcription errors. They are not interval certificates or substitutes
for those proofs.

% END sections/00_scope.tex

% BEGIN sections/01_dougall.tex
\section{Centered Dougall subtraction and every integer resonance}
\label{dg:section}

\subsection{The summation family and its natural center}
Fix positive real parameters $a,b,c$ such that
\begin{equation}\label{dg:domain}
 d_0=1+a-b-c>0,\qquad \eta=\frac a2,\qquad x_n=n+\eta.
\end{equation}
The restriction $d_0>0$ supplies a common domain of ordinary convergence
from which all continuations below are proved. It still includes the
quartic-binomial examples and all their nonpositive-integer resonances.
For $\Re u<d_0$ define
\begin{equation}\label{dg:rn}
 r_n(u)=x_n\,
 \frac{\Gamma(n+a)\Gamma(n+b)\Gamma(n+c)\Gamma(n+d_0-u)}
 {\Gamma(n+1)\Gamma(n+1+a-b)\Gamma(n+1+a-c)
  \Gamma(n+b+c+u)}.
\end{equation}
All reciprocal-Gamma factors are entire functions. In particular,
$r_n$ is holomorphic when its last denominator argument passes through
a nonpositive integer. No numerator Gamma pole lies in the stated
half-plane.

The Rogers--Dougall summation \cite[Eq.~16.4.9]{dlmf-hyper} gives
\begin{equation}\label{dg:sum}
 \sum_{n=0}^{\infty}r_n(u)=G(u):=
 \frac{\Gamma(b)\Gamma(c)\Gamma(d_0-u)\Gamma(u)}
 {2\Gamma(d_0)\Gamma(b+u)\Gamma(c+u)},
 \qquad 0<\Re u<d_0.
\end{equation}
To check the normalization, the upper parameters of the usual
${}_5F_4$ are $a,1+a/2,b,c,d_0-u$ and its lower parameters are
$a/2,1+a-b,1+a-c,b+c+u$. The ratio
$(1+a/2)_n/(a/2)_n$ is $(n+a/2)/(a/2)$, which accounts for the
factor $1/2$ in \eqref{dg:sum}. The conventional hypergeometric excess
is $2u$, whereas the Gamma factor in the sum is $\Gamma(u)$.
Keeping that distinction is essential at every resonance.

\subsection{Even inverse powers and finite coefficient data}
Write
\begin{equation}\label{dg:alpha}
 (\alpha_1,\alpha_2,\alpha_3,\alpha_4)
 =\left(\frac a2,b-\frac a2,c-\frac a2,d_0-u-\frac a2\right).
\end{equation}
At the center $x_n=n+a/2$, every denominator shift is the reflection
$1-\alpha_i$ of its numerator shift:
\[
 r_n(u)=x_n\prod_{i=1}^4
       \frac{\Gamma(x_n+\alpha_i)}{\Gamma(x_n+1-\alpha_i)}.
\]
For Bernoulli polynomials with $B_1(z)=z-1/2$, define
\begin{align}
 h_k(u)&=-\frac{1}{k(2k+1)}
                \sum_{i=1}^4B_{2k+1}(\alpha_i),\qquad k\ge1,
                  \label{dg:hk}\\
 C_0(u)&=1,\qquad
 jC_j(u)=\sum_{k=1}^j k h_k(u)C_{j-k}(u),\qquad j\ge1.
                  \label{dg:recurrence}
\end{align}
These are polynomials with rational coefficients in $a,b,c,u$; each
prescribed order requires only finitely many polynomial operations.

\begin{lemma}[Centered finite expansion]\label{dg:asymptotic}
For every $K\ge1$,
\begin{equation}\label{dg:expansion}
 r_n(u)=x_n^{-1-2u}
 \left\{\sum_{j=0}^{K-1}C_j(u)x_n^{-2j}+O(x_n^{-2K})\right\}.
\end{equation}
The remainder is locally uniform in $\Re u<d_0$ and in the positive
parameter domain \eqref{dg:domain}. Every fixed-order derivative has
the corresponding differentiated expansion, with at most a fixed
power of $\log x_n$.
\end{lemma}
\begin{proof}
The finite Stirling--Bernoulli expansion for the logarithm of a Gamma
ratio \cite[Section 5.11]{dlmf-gamma} contains the coefficient
\[
 \frac{(-1)^{k+1}}{k(k+1)}
 \{B_{k+1}(\alpha_i)-B_{k+1}(1-\alpha_i)\}x_n^{-k}.
\]
Bernoulli reflection makes this zero for odd $k$. For $k=2\ell$ its
sum over $i$ is $h_\ell(u)$. The leading power is $x_n^{-1-2u}$.
Exponentiation gives \eqref{dg:expansion}, and differentiation of the
formal exponential in the inverse-square variable gives
\eqref{dg:recurrence}. This argument uses only a finite asymptotic
expansion with its remainder; it does not sum a divergent asymptotic
series. Uniformity follows from the uniform Stirling remainder on
compact sets of shifts. To justify parameter derivatives, enlarge a
given compact set slightly and apply Cauchy's formula to the remainder.
The finitely many small indices are handled by the entire
reciprocal-Gamma factors.
\end{proof}

\begin{theorem}[Centered Dougall--Hurwitz subtraction]\label{dg:master}
For $K\ge1$, the series
\begin{equation}\label{dg:RK}
 R_K(u)=\sum_{n=0}^{\infty}\left\{r_n(u)-
    \sum_{j=0}^{K-1}C_j(u)x_n^{-1-2u-2j}\right\}
\end{equation}
converges absolutely and locally normally on $-K<\Re u<d_0$ and satisfies
\begin{equation}\label{dg:mastereq}
 \boxed{R_K(u)=G(u)-
       \sum_{j=0}^{K-1}C_j(u)\zeta(1+2u+2j,\eta).}
\end{equation}
All poles of the combined right side in this strip are removable.
Every fixed spectral or parameter derivative can be taken in the
single bracketed series \eqref{dg:RK}.
\end{theorem}
\begin{proof}
On a compact subset where $\Re u\ge-K+\delta$, the bracket is
$O(n^{-1-2\delta})$ by Lemma~\ref{dg:asymptotic}. Derivatives add only
fixed logarithmic powers, so the same argument gives summable
majorants. The Weierstrass theorem proves holomorphy and the derivative
assertion. In $0<\Re u<d_0$, summation of the separate convergent
series using \eqref{dg:sum} and the Hurwitz Dirichlet series gives
\eqref{dg:mastereq}. The identity theorem extends it to the stated
strip. Its left side proves removability of the combined poles.
\end{proof}

\begin{remark}[The evaluated object]\label{dg:ordinary}
At $u=-N$, the two components in braces usually have divergent sums
separately. Equation~\eqref{dg:mastereq} evaluates their one absolutely
convergent difference. The reference $\eta=a/2$ and the number of
subtractions are part of the definition. Increasing that number obeys
the exact relation
\[
 R_{K+1}(u)=R_K(u)-C_K(u)\zeta(1+2u+2K,\eta),
 \qquad -K<\Re u<d_0.
\]
\end{remark}

\subsection{The coefficient differential law}
The even coefficients have an exact parameter identity that considerably
simplifies their resonant constant terms. Every derivative in the
following operator respects $d_0=1+a-b-c$:
\[
 V=2\partial_a+\partial_b+\partial_c-\partial_u.
\]

\begin{lemma}[Centered coefficient transport]\label{dg:differential}
For every $j\ge1$,
\begin{equation}\label{dg:differentiallaw}
 \boxed{VC_j(u)=2\sum_{k=1}^j
                      \zeta(1-2k,\eta)C_{j-k}(u).}
\end{equation}
This is a polynomial identity, since
$\zeta(1-2k,\eta)=-B_{2k}(\eta)/(2k)$.
\end{lemma}
\begin{proof}
Temporarily regard $n$ as a continuous positive variable. Direct
logarithmic differentiation of \eqref{dg:rn} gives
\begin{equation}\label{dg:termtransport}
 Vr_n(u)=\partial_n r_n(u)
                 +\{\psi(n+a)+\psi(n+1)\}r_n(u).
\end{equation}
Here $Vx_n=\partial_nx_n=1$. The reflected digamma pair has expansion
\[
 \psi(x_n+\eta)+\psi(x_n+1-\eta)
 \sim2\log x_n+2\sum_{k\ge1}\zeta(1-2k,\eta)x_n^{-2k}.
\]
This again follows from Bernoulli reflection in the differentiated
Stirling formula. Since $Vu=-1$, the explicit exponent on the left
contributes $2\log x_n$, which cancels the displayed logarithm.
The derivatives of $x_n$ cancel between $V$ and $\partial_n$.
Comparison of each inverse-square coefficient proves
\eqref{dg:differentiallaw}. Uniqueness of finite asymptotic
coefficients justifies this comparison.
\end{proof}

\subsection{An explicit residue polynomial and a short closed sum}
For $N\ge0$ set
\begin{equation}\label{dg:rho}
 \rho_N=\frac{(-1)^N(d_0)_N(1-b)_N(1-c)_N}{N!}.
\end{equation}

\begin{proposition}[Two finite resonance identities]\label{dg:residues}
For every $N\ge0$,
\begin{align}
 C_N(-N)&=\rho_N,\label{dg:rhoidentity}\\
 \frac12C_N'(-N)+
 \sum_{j=0}^{N-1}C_j(-N)\zeta(1-2N+2j,\eta)
 &=\left(\partial_a+\tfrac12\partial_b+\tfrac12\partial_c\right)\rho_N.
 \label{dg:Bernoulliidentity}
\end{align}
The sum is empty when $N=0$. The formulas hold also at all zeros of
$\rho_N$.
\end{proposition}
\begin{proof}
Use $K=N+1$ in Theorem~\ref{dg:master}. The only Hurwitz pole at
$u=-N$ has index $j=N$ and residue $C_N(-N)/2$. The residue of
$G(u)$ is $\rho_N/2$, using the Gamma residue and shift formulas.
Since the bracketed sum is holomorphic, their residues agree. This
proves \eqref{dg:rhoidentity} generically and then identically as a
polynomial statement. Evaluate \eqref{dg:differentiallaw} at $u=-N$,
use \eqref{dg:rhoidentity} for its parameter derivatives, and divide
by two. This gives \eqref{dg:Bernoulliidentity}; the case $N=0$ is
immediate from $C_0=\rho_0=1$.
\end{proof}

\begin{theorem}[Every nonpositive-integer resonance]\label{dg:closed}
Under \eqref{dg:domain}, for every integer $N\ge0$,
\begin{align}
 &\sum_{n=0}^{\infty}\left\{r_n(-N)-
             \sum_{j=0}^{N}C_j(-N)x_n^{2N-1-2j}\right\}\notag\\
 &\qquad=\boxed{\rho_N\left[
   \psi(\eta)+\frac{\psi(N+1)-\psi(b)-\psi(c)-\psi(d_0+N)}2
                    \right].}\label{dg:closedeq}
\end{align}
The series is absolutely convergent, including at positive integer
parameters where $\rho_N=0$.
\end{theorem}
\begin{proof}
First suppose that the shifted Gamma arguments in the following
Laurent expansion are regular. With $u=-N+t$,
\[
 G(-N+t)=\frac{\rho_N}{2t}
 +\frac{\rho_N}{2}
   \{\psi(N+1)-\psi(d_0+N)-\psi(b-N)-\psi(c-N)\}+O(t).
\]
The colliding Hurwitz term has constant coefficient
$C_N'(-N)/2-\rho_N\psi(\eta)$. Consequently
\eqref{dg:mastereq}, at $K=N+1$, gives the claimed value except for
the finite expression on the left of \eqref{dg:Bernoulliidentity}.
The operator $W=\partial_a+(\partial_b+\partial_c)/2$ fixes $d_0$, and
\[
 W\rho_N=\frac{\rho_N}{2}
   \{\psi(b)-\psi(b-N)+\psi(c)-\psi(c-N)\}.
\]
Subtraction cancels both shifted digammas and proves
\eqref{dg:closedeq}. Both the normally convergent sum and the final
right side are locally holomorphic in the positive parameter domain.
They therefore agree across zeros of $\rho_N$, without division by
that polynomial at its zeros.
\end{proof}

At the resonant point write $d=d_0+N$. Then
$a=b+c+d-1-N$ and
\[
 \rho_N=\frac{(1-b)_N(1-c)_N(1-d)_N}{N!},
\]
which displays the symmetry of the answer in $b,c,d$. Our proof and
theorem retain $d_0>0$; no continuation into an unspecified larger
parameter domain is needed.

\subsection{Half-integer excess and parity cancellation}
Negative odd values of the conventional excess $2u$ behave differently.
At $u=-M-1/2$, the minimal centered expansion subtracts the nonnegative
even powers $x_n^{2M},x_n^{2M-2},\ldots,1$. There is no $x_n^{-1}$
coefficient to collide with a Hurwitz pole.

\begin{corollary}[Negative odd hypergeometric excess]\label{dg:half}
For every integer $M\ge0$,
\begin{equation}\label{dg:halfeq}
 R_{M+1}(-M-\tfrac12)
 =G(-M-\tfrac12)-\sum_{j=0}^{M}
       C_j(-M-\tfrac12)\zeta(-2M+2j,\eta).
\end{equation}
All terms on the right are individually regular. In particular,
\begin{align}
 \sum_{n=0}^{\infty}\{r_n(-\tfrac12)-1\}
 =\frac{a-1}{2}
 -\frac{\sqrt\pi\,\Gamma(b)\Gamma(c)\Gamma(d_0+\tfrac12)}
 {\Gamma(d_0)\Gamma(b-\tfrac12)\Gamma(c-\tfrac12)}.
 \label{dg:halfsimple}
\end{align}
\end{corollary}
\begin{proof}
Theorem~\ref{dg:master} applies with $K=M+1$ since
$-M-1/2>-M-1$. Its Hurwitz arguments are exactly the nonpositive
even integers in \eqref{dg:halfeq}. The case $M=0$ uses
$\Gamma(-1/2)=-2\sqrt\pi$ and $\zeta(0,\eta)=1/2-\eta$.
\end{proof}

For example, $a=b=c=3/4$ and $d_0=1/4$ give
\begin{equation}\label{dg:quarteridentity}
 \boxed{\sum_{n=0}^{\infty}\left\{
 (n+\tfrac38)\left[\frac{\Gamma(n+\tfrac34)}{\Gamma(n+1)}\right]^4-1
 \right\}=-\frac18-\sqrt\pi
     \left[\frac{\Gamma(\tfrac34)}{\Gamma(\tfrac14)}\right]^3.}
\end{equation}

A complementary specialization has a vanishing Gamma endpoint.
For $a=b=c=1/2$ and $u=-1/2$, put
\[
 b_n=\pi\frac{\binom{2n}{n}^2}{16^n},\qquad
 v_n=(n+\tfrac14)b_n.
\]
The two reciprocal-Gamma factors at zero give $G(-1/2)=0$, so
\begin{equation}\label{dg:binomialsquare}
 \sum_{n\ge0}(v_n-1)=-\frac14.
\end{equation}
There is also a direct check: the recurrence
$(n+1)^2b_{n+1}=(n+1/2)^2b_n$ gives
$\sum_{n=0}^{M-1}v_n=M^2b_M$. Stirling's formula yields
$M^2b_M=M-1/4+O(M^{-1})$, proving \eqref{dg:binomialsquare}
independently of endpoint continuation.

At $u=-M-1/2$, the same parameters give
\[
 r_n(-M-\tfrac12)=v_n\prod_{r=0}^{M-1}
           \{(n+\tfrac14)^2-(r+\tfrac34)^2\},\qquad
 G(-M-\tfrac12)=0.
\]
Their minimally subtracted sums are rational numbers, because
\eqref{dg:halfeq} involves rational polynomial coefficients and
Hurwitz values at nonpositive integers and shift $1/4$. For example,
with $x_n=n+1/4$,
\begin{equation}\label{dg:binomialsquareN1}
 \sum_{n\ge0}\left\{(x_n^2-\tfrac9{16})v_n-x_n^2+\tfrac{19}{32}\right\}
 =\frac{21}{128}.
\end{equation}
The next section develops an integer-resonance family whose values
belong to $\Q+\pi\Q$.

Parity persists at an arbitrary reference shift $v>0$, although the
individual odd coefficients then need not vanish. If
\[
 r_n(u)\sim(n+v)^{-1-2u}\sum_{j\ge0}A_j(u;v)(n+v)^{-j},
\]
the binomial theorem gives the exact finite transport formula
\begin{equation}\label{dg:reference}
 A_j(u;v)=\sum_{\ell=0}^{\lfloor j/2\rfloor}
 C_\ell(u)\frac{(1+2u+2\ell)_{j-2\ell}}{(j-2\ell)!}
                       (v-\eta)^{j-2\ell}.
\end{equation}
At the potentially colliding index it implies
\begin{equation}\label{dg:parityresidue}
 A_{2N}(-N;v)=\rho_N,\qquad
 A_{2N+1}(-N-\tfrac12;v)=0.
\end{equation}
For the odd index every rising factorial in \eqref{dg:reference}
contains a zero; for the even index only the term $\ell=N$ survives.
Thus a generic all-powers subtraction scheme must exhibit the same
residue cancellation. Derivatives of a coefficient vanishing at the
resonant point need not vanish, a distinction used next.

% END sections/01_dougall.tex

% BEGIN sections/02_dougall_jets.tex
\section{Complete Dougall jets and quartic harmonic identities}
\label{dg:jetssection}

\subsection{A holomorphic Gamma factor at the collision}
Fix $N\ge0$. The factorization
\begin{equation}\label{dg:EN}
 E_N(t)=
 \frac{\Gamma(b)\Gamma(c)\Gamma(d_0+N-t)\Gamma(1+t)}
 {\Gamma(d_0)\Gamma(b+t)\Gamma(c+t)}
 \prod_{r=1}^{N}\frac{(b+t-r)(c+t-r)}{t-r}
\end{equation}
follows from the Gamma functional equation and satisfies
\begin{equation}\label{dg:ENidentity}
 E_N(t)=2tG(-N+t),\qquad E_N(0)=\rho_N.
\end{equation}
Every denominator in the finite product is nonzero near $t=0$.
Positive integer $b$ or $c$ may create numerator zeros; those zeros
remain in the polynomial and are differentiated with it. Consequently
\eqref{dg:EN} computes all Taylor coefficients without a singular
Gamma quotient or a division by $\rho_N$.

Put
\begin{equation}\label{dg:jetcoeffs}
 e_q=[t^q]E_N(t),\qquad c_{j,q}=[t^q]C_j(-N+t).
\end{equation}
The coefficients $c_{j,q}$ are finite polynomial data. The $e_q$ are
obtained from Bell polynomials in Gamma logarithmic derivatives at
positive arguments, followed by the finite product in \eqref{dg:EN}.

\begin{theorem}[All resonant Stieltjes coefficients]\label{dg:alljets}
For $K>N$ and $m\ge0$,
\begin{align}
 [t^m]R_K(-N+t)
 ={}&\frac{e_{m+1}-c_{N,m+1}}2
 -\sum_{h=0}^{m}c_{N,m-h}\frac{(-2)^h\gamma_h(\eta)}{h!}
 \notag\\
 &-\sum_{\substack{0\le j<K\\j\ne N}}\sum_{q=0}^{m}
 c_{j,q}\frac{2^{m-q}\zeta^{(m-q)}(1-2N+2j,\eta)}{(m-q)!}.
 \label{dg:alljetseq}
\end{align}
The left side equals the absolutely convergent ordinary series
\begin{equation}\label{dg:jetseries}
 \sum_{n=0}^{\infty}[t^m]\left\{r_n(-N+t)
 -\sum_{j=0}^{K-1}C_j(-N+t)x_n^{2N-1-2j}e^{-2t\log x_n}\right\}.
\end{equation}
At $K=N+1$, the zeta derivatives in the last line of
\eqref{dg:alljetseq} have only the negative odd arguments
$1-2N,3-2N,\ldots,-1$.
\end{theorem}
\begin{proof}
Normal convergence in Theorem~\ref{dg:master} permits coefficient
extraction in its bracketed series. The Gamma term is $E_N(t)/(2t)$.
The singular Hurwitz factor is
\[
 \zeta(1+2t,\eta)=\frac1{2t}
          +\sum_{h\ge0}\frac{(-2)^h\gamma_h(\eta)}{h!}t^h.
\]
Multiplication by $C_N(-N+t)$ contributes
$c_{N,m+1}/2$ and the first finite sum in \eqref{dg:alljetseq}.
Every other Hurwitz factor is regular and expands by the ordinary
Cauchy product, giving the second finite sum. The coefficient of
$t^{-1}$ cancels by \eqref{dg:rhoidentity}. This establishes the
formula without a limiting prescription at a pole.
\end{proof}

\begin{corollary}[Finite special-function closure]\label{dg:closure}
At fixed $N,m$ and $K=N+1$, the resonant Taylor coefficient has a
finite expression in Gamma and polygamma data at $b,c,d_0+N,1$,
generalized Stieltjes constants $\gamma_h(\eta)$ with $0\le h\le m$,
and spectral derivatives $\zeta^{(q)}(1-2r,\eta)$ with
$1\le r\le N$ and $0\le q\le m$. Coefficients from finite Gamma
shifts are polynomial or rational with nonsingular denominators.
\end{corollary}
This is a representation theorem. It asserts neither algebraic
independence nor a minimal special-value basis. Negative-order
Hurwitz derivatives can be changed to integrated-polygamma or
Barnes-type normalizations, subject to the normalization conventions
already emphasized in the incoming reports \cite{incoming-gh,incoming-sh}.

At a half-integer point $u_0=-M-1/2$ there is no singular factor.
Writing $\widehat c_{j,q}=[t^q]C_j(u_0+t)$, all its jets are simply
\begin{align}
 [t^m]R_{M+1}(u_0+t)
 ={}&[t^m]G(u_0+t)\notag\\
 &-\sum_{j=0}^{M}\sum_{q=0}^{m}\widehat c_{j,q}
       \frac{2^{m-q}\zeta^{(m-q)}(-2M+2j,\eta)}{(m-q)!}.
 \label{dg:halfjets}
\end{align}
Only nonpositive even spectral arguments occur; no generalized
Stieltjes constant is required in this centered normalization.

\subsection{An explicit Bell identity at zero excess}
Define, for $k\ge1$,
\begin{align}
 L_k={}&(-1)^k\psi^{(k-1)}(d_0)+\psi^{(k-1)}(1)
                    -\psi^{(k-1)}(b)-\psi^{(k-1)}(c),
                    \label{dg:Lk}\\
 T_k(n)={}&(-1)^k\psi^{(k-1)}(n+d_0)
                    -\psi^{(k-1)}(n+b+c).
                    \label{dg:Tk}
\end{align}
The $L_k$ are the logarithmic derivatives of $E_0(t)$, while $T_k(n)$
are those of $r_n(t)$ at $t=0$. With the Bell normalization stated
in Section~\ref{sec:scope}, Theorem~\ref{dg:alljets} gives
\begin{align}
 &\sum_{n=0}^{\infty}\left\{
 r_n(0)Y_m(T_1(n),\ldots,T_m(n))
              -\frac{(-2\log x_n)^m}{x_n}\right\}\notag\\
 &\hspace{18mm}=
 \boxed{\frac{Y_{m+1}(L_1,\ldots,L_{m+1})}{2(m+1)}
              -(-2)^m\gamma_m(\eta).}\label{dg:Bellzero}
\end{align}
Here $Y_0=1$. The formula is for derivatives; the multiplication by
$m!$ relative to \eqref{dg:alljetseq} explains its factor $1/(m+1)$.
For instance,
\begin{align}
 \sum_{n\ge0}\left\{r_n(0)T_1(n)+\frac{2\log x_n}{x_n}\right\}
 &=\frac{L_1^2+L_2}{4}+2\gamma_1(\eta),\label{dg:Bellone}\\
 \sum_{n\ge0}\left\{r_n(0)(T_1(n)^2+T_2(n))
                    -\frac{4\log^2x_n}{x_n}\right\}
 &=\frac{L_1^3+3L_1L_2+L_3}{6}-4\gamma_2(\eta).
 \label{dg:Belltwo}
\end{align}
The paired logarithmic counterterms are essential for convergence.

\subsection{Finite polynomial shifts preserve every moving zero}
There is a convenient exact identity at the summand level:
\begin{equation}\label{dg:polynomialshift}
 \boxed{r_n(-N+t)=r_n(t)
             (n+d_0-t)_N(n+b+c-N+t)_N.}
\end{equation}
It follows from the Gamma functional equation. The base $r_n(t)$ has
positive Gamma arguments at $t=0$; the finite polynomial contains
every shifted reciprocal-Gamma zero. Thus all resonant summand jets
are finite convolutions of its polynomial coefficients with the Bell
coefficients in \eqref{dg:Tk}. This provides a direct implementation
for all orders without singular harmonic expressions at small $n$.

The distinction is substantive. A term that is zero after substituting
$u=-N$ can have nonzero first and higher derivatives. Deleting such
terms before constructing the deformation changes the value of the
jet. The following family supplies a simple exact example.

\subsection{Quartic central-binomial moments at all resonances}
Set $a=b=c=d_0=1/2$ and put
\begin{align}
 x_n&=n+\tfrac14,\notag\\
 w_n&=x_n\left[\frac{\Gamma(n+\tfrac12)}{\Gamma(n+1)}\right]^4
      =\pi^2x_n\frac{\binom{2n}{n}^4}{256^n},
                   \label{dg:wn}\\
 P_N(X)&=\prod_{r=0}^{N-1}\{X-(r+\tfrac14)^2\}.
                   \label{dg:PN}
\end{align}
The finite Gamma shift gives
\begin{equation}\label{dg:quarticterm}
 r_n(-N)=w_nP_N(x_n^2),\qquad
 \rho_N=\frac{(-1)^N(\tfrac12)_N^3}{N!}.
\end{equation}
Using
\[
 \psi(\tfrac14)=-\gamma-3\log2-\frac\pi2,
 \qquad
 \psi(N+\tfrac12)=-\gamma-2\log2+2H_{2N}-H_N,
\]
Theorem~\ref{dg:closed} becomes the all-resonance identity
\begin{align}
 &\sum_{n=0}^{\infty}\left\{w_nP_N(x_n^2)
            -\sum_{j=0}^{N}C_j(-N)x_n^{2N-1-2j}\right\}\notag\\
 &\qquad=\boxed{\frac{(-1)^N(\tfrac12)_N^3}{N!}
                  \left(H_N-H_{2N}-\frac\pi2\right).}
                  \label{dg:quarticall}
\end{align}
All counterterms still follow from the finite recurrence
\eqref{dg:hk}--\eqref{dg:recurrence}; for example,
\begin{equation}\label{dg:quarticC1}
 C_1(u)=\frac{u^3}{3}+\frac{u^2}{4}-\frac{u}{48}-\frac1{16}.
\end{equation}
They can also be built from one universal base sequence. If $\sigma_k$ denotes
the elementary symmetric polynomial in
$(1/4)^2,(5/4)^2,\ldots,(N-3/4)^2$, with $\sigma_0=1$, then
\begin{equation}\label{dg:quarticcoefficienttransport}
 C_j(-N)=\sum_{k=0}^{\min(j,N)}(-1)^k\sigma_k C_{j-k}(0).
\end{equation}
This follows by multiplying the polynomial $P_N(x_n^2)$ into the
base expansion of $w_n$. Its first base coefficients are
$C_0(0)=1$, $C_1(0)=-1/16$, $C_2(0)=3/256$, and
$C_3(0)=-23/4096$. Thus successive moment subtractions do not require
an independent asymptotic calculation.
The first three cases are
\begin{align}
 \sum_{n\ge0}\left\{w_n-\frac1{x_n}\right\}
       &=-\frac\pi2,\label{dg:quarticN0}\\
 \sum_{n\ge0}\left\{n(n+\tfrac12)w_n-x_n+\frac1{8x_n}\right\}
       &=\frac{\pi+1}{16},\label{dg:quarticN1}\\
 \sum_{n\ge0}\left\{n(n-1)(n+\tfrac12)(n+\tfrac32)w_n
      -x_n^3+\frac{27}{16}x_n-\frac{27}{128x_n}\right\}
       &=-\frac{63+54\pi}{512}.\label{dg:quarticN2}
\end{align}
The zero-excess case \eqref{dg:quarticN0} is a recovered classical
Ramanujan--Dougall constant, explicitly included in B\"uhring's
partial-sum analysis \cite[Section 4, Corollaries 1--2]{buhring}.
The uniform resonant mechanism and its complete jets are the
continuation supplied here. The displayed higher examples are proved
specializations, without a claim that each has no earlier appearance.

\subsection{Generalized harmonic coefficients of every weight}
For the same half-parameter family, write
\[
 g=\gamma+\log2,\qquad g_3=\gamma+3\log2.
\]
Then $T_1(n)=2(g-H_{2n})$, $L_1=2g_3$, and for $k\ge2$,
\begin{align}
 T_k(n)=(k-1)!\big\{&
 (2^k-1-(-1)^k)\zeta(k)-2^kH_{2n}^{(k)}
                  +(1+(-1)^k)H_n^{(k)}\big\},
                    \label{dg:harmonicTk}\\
 L_k=(k-1)!\big\{&(1-2(-1)^k)(2^k-1)+(-1)^k\big\}\zeta(k).
                    \label{dg:harmonicLk}
\end{align}
These follow by substituting the standard half-integer polygamma
values into \eqref{dg:Lk}--\eqref{dg:Tk}. In particular,
$L_2=-2\zeta(2)$, $L_3=40\zeta(3)$, and $L_4=-84\zeta(4)$.
Together with \eqref{dg:Bellzero}, they give an explicit all-weight
family of generalized harmonic identities. Its first nonconstant
member is
\begin{equation}\label{dg:quarticharmonic1}
 \boxed{\sum_{n\ge0}\left\{w_n(g-H_{2n})
                         +\frac{\log x_n}{x_n}\right\}
       =\frac{g_3^2}{2}-\frac{\zeta(2)}4+\gamma_1(\tfrac14).}
\end{equation}

At higher resonance, the stronger finite identity
\begin{equation}\label{dg:quarticshift}
 r_n(-N+t)=r_n(t)\prod_{r=0}^{N-1}
                   \{x_n^2-(r+\tfrac14-t)^2\}
\end{equation}
supplies all harmonic summand coefficients by finite convolution.
For $n<N$, put $k=N-n-1$. The first revived term is
\begin{equation}\label{dg:quarticrevival}
 [t]r_n(-N+t)=w_n(n+\tfrac12)_N n!(-1)^k k!.
\end{equation}
Indeed, the simple zero of $1/\Gamma(-k+t)$ has leading coefficient
$(-1)^kk!t$; the other factors are regular. In particular,
\begin{equation}\label{dg:revivalpi}
 r_0(-1)=0,\qquad [t]r_0(-1+t)=\frac{\pi^2}{8}.
\end{equation}

For completeness, the full next-resonance first derivative can be
written without any exceptional small index. For every $n\ge0$,
\[
 [t]r_n(-1+t)=w_n\{2n(n+\tfrac12)(g-H_{2n})+\tfrac12\}.
\]
Also
\[
 C_1(-1+t)=-\frac18+\frac{23}{48}t-\frac34t^2+\frac13t^3.
\]
Substituting these coefficients in \eqref{dg:alljetseq} yields
\begin{align}
 \sum_{n\ge0}\Bigg\{&
 w_n\big[2n(n+\tfrac12)(g-H_{2n})+\tfrac12\big]
 +2x_n\log x_n
 -\frac{\tfrac14\log x_n+\tfrac{23}{48}}{x_n}\Bigg\}
 \notag\\
 ={}&-\frac{g_3^2}{8}+\frac{7g_3}{48}+\frac{\zeta(2)}{16}
 -\frac1{16}-\frac{23\pi}{96}
 -\frac{\gamma_1(\tfrac14)}4-2\zeta^{(1)}(-1,\tfrac14).
 \label{dg:quarticN1jet}
\end{align}
One way to verify the constant directly is to use
$E_1(t)=E_0(t)(t-1/2)^3/(1-t)$ in \eqref{dg:EN}, and
\[
 [t^2]E_1(t)=-\frac{L_1^2}{16}
                +\frac{5L_1}{8}-\frac{L_2}{16}-\frac78.
\]
The formula combines a generalized harmonic summand, a Stieltjes
constant, a negative-odd Hurwitz derivative, and the revived finite
head. Deleting the $n=0$ summand after first substituting $u=-1$
would change the first derivative by exactly $\pi^2/8$.

\subsection{Squared-binomial harmonic jets and log-Gamma}
The squared-binomial half-resonance family in
\eqref{dg:binomialsquare} has an additional cancellation. At
$a=b=c=1/2$,
\begin{equation}\label{dg:halfdoublezero}
 G(-\tfrac12+t)=\frac{\sqrt\pi}{2}
       \frac{\Gamma(1-t)\Gamma(-\tfrac12+t)}{\Gamma(t)^2}
       =-\pi t^2+O(t^3).
\end{equation}
Thus its constant and first Taylor coefficients vanish. Recall
$v_n=\pi(n+1/4)\binom{2n}{n}^2/16^n$, $x_n=n+1/4$, and
$g=\gamma+\log2$. The logarithmic derivatives of
$r_n(-1/2+t)$ at $t=0$ are
\[
 2(g-H_{2n}),\qquad
 -2\zeta(2)+4H_{2n}^{(2)}-2H_n^{(2)}
\]
at orders one and two. Applying \eqref{dg:halfjets} with $M=0$,
and using Lerch's identity for $\zeta^{(1)}(0,1/4)$
\cite[Eq.~25.11.18]{dlmf-hurwitz}, gives
\begin{equation}\label{dg:squaredharmonic1}
 \boxed{\sum_{n\ge0}\{v_n(g-H_{2n})+\log x_n\}
 =\frac12\log(2\pi)-\log\Gamma(\tfrac14).}
\end{equation}
The second derivative gives
\begin{align}
 \sum_{n\ge0}\Big\{v_n\big[(g-H_{2n})^2
       -\tfrac12\zeta(2)+H_{2n}^{(2)}-\tfrac12H_n^{(2)}\big]
                         -\log^2x_n\Big\}
 =-\frac\pi2-\zeta^{(2)}(0,\tfrac14).
 \label{dg:squaredharmonic2}
\end{align}
Both are single absolutely convergent bracketed sums. Equations
\eqref{dg:halfjets} and \eqref{dg:halfdoublezero} similarly determine
every higher harmonic jet. Here the pole-free half-resonance produces
Hurwitz derivatives at zero, while the integer-resonance family
\eqref{dg:Bellzero} produces generalized Stieltjes constants at the
pole. The distinction follows from the exact centered spectrum.

% END sections/02_dougall_jets.tex

% BEGIN sections/03_polylog_primitives.tex
\section{A hypergeometric primitive and finite polylogarithmic transport}
\label{dp:section}

\subsection{A coefficientwise regularized \texorpdfstring{${}_4F_3$}{4F3}}
The centered prefactor has a useful consequence for antiderivatives.
For the same parameter domain, set
\begin{equation}\label{dp:Q}
 q_n(u)=\frac{r_n(u)}{n+\eta},\qquad
 Q_u(z)=\sum_{n=0}^{\infty}q_n(u)z^n,\qquad |z|<1.
\end{equation}
Generically this is the Gamma-normalized hypergeometric function
\begin{align}
 Q_u(z)={}&
 \frac{\Gamma(a)\Gamma(b)\Gamma(c)\Gamma(d_0-u)}
 {\Gamma(1+a-b)\Gamma(1+a-c)\Gamma(b+c+u)}\notag\\
 &\times{}_4F_3\!\left(\begin{matrix}
 a,b,c,d_0-u\\1+a-b,1+a-c,b+c+u
 \end{matrix};z\right).
 \label{dp:4F3}
\end{align}
At a nonpositive lower parameter, the product in \eqref{dp:4F3} is
defined by the coefficient series \eqref{dp:Q}, or equivalently by
regularizing the lower Gamma factors before evaluation. The ordinary
${}_4F_3$ and its prefactor must not be separately evaluated as an
undefined expression times zero. The definition \eqref{dp:Q} retains
the zero summands and all their parameter derivatives.

Let $\theta=z\,\mathrm d/\mathrm dz$ and use the Lerch transcendent
\cite[Section 25.14]{dlmf-lerch}
\[
 \Phi(z,s,\eta)=\sum_{n=0}^{\infty}\frac{z^n}{(n+\eta)^s},
 \qquad |z|<1.
\]
The subtracted generating function is
\begin{equation}\label{dp:F}
 F_{K,u}(z)=(\theta+\eta)Q_u(z)
        -\sum_{j=0}^{K-1}C_j(u)\Phi(z,1+2u+2j,\eta).
\end{equation}
It is exactly the generating function of the brackets in
\eqref{dg:RK}. For $-K<\Re u<d_0$, their absolute summability gives
\begin{equation}\label{dp:endpoint}
 \lim_{z\uparrow1}F_{K,u}(z)=R_K(u).
\end{equation}
The limit and every fixed spectral derivative are locally uniform on
compact subsets of this strip. The identity follows from the
summable majorant in Theorem~\ref{dg:master}, not merely from an
unqualified appeal to boundary continuation.

\subsection{An exact weighted primitive}
\begin{theorem}[The ${}_5F_4$-to-${}_4F_3$ primitive]\label{dp:primitive}
For $|z|<1$ on a fixed branch of $z^\eta$, the normalized radial
primitive of \eqref{dp:F} is
\begin{align}
 \int_0^z t^{\eta-1}F_{K,u}(t)\dd t
  =z^\eta\left\{Q_u(z)-\sum_{j=0}^{K-1}
            C_j(u)\Phi(z,2+2u+2j,\eta)\right\}.
 \label{dp:primitiveeq}
\end{align}
It vanishes at zero. Every fixed spectral derivative and every
parameter derivative within the domain may be taken through this
identity, including the dependence of the weight on $\eta=a/2$.
\end{theorem}
\begin{proof}
The defining series converge normally for $|z|<1$ on compact parameter
sets. Integrating the coefficient of $z^n$ uses
\[
 \int_0^z t^{n+\eta-1}\dd t=\frac{z^{n+\eta}}{n+\eta}.
\]
Thus the coefficient $r_n(u)=(n+\eta)q_n(u)$ loses its centered
factor, while the order of every Lerch coefficient increases by one.
This proves \eqref{dp:primitiveeq}. The condition $\eta>0$ makes
the integrand locally integrable at zero and the displayed primitive
zero there. On slightly enlarged compact parameter sets, additional
logarithmic factors from differentiating $t^{\eta-1}$ are still
integrable. Normal convergence and dominated integration therefore
give the derivative assertions.
\end{proof}

More generally, division of coefficients by $(n+\eta)^\ell$ gives
successive applications of the normalized inverse of $\theta+\eta$:
\begin{align}
 F_{K,u}^{\langle\ell\rangle}(z)
 ={}&\sum_{n=0}^{\infty}
       \frac{r_n(u)z^n}{(n+\eta)^\ell}\notag\\
 &-\sum_{j=0}^{K-1}C_j(u)
            \Phi(z,1+2u+2j+\ell,\eta),\qquad \ell\ge0,
 \label{dp:ladder}\\
 (\theta+\eta)F_{K,u}^{\langle\ell+1\rangle}
 ={}&F_{K,u}^{\langle\ell\rangle}.
 \label{dp:ladderrec}
\end{align}
Here $F_{K,u}^{\langle0\rangle}=F_{K,u}$ and the first step is the
lower-order hypergeometric expression in \eqref{dp:primitiveeq}.
All these identities are coefficientwise regular at moving lower
Gamma zeros.

\subsection{Rational shifts give a finite colored-polylogarithm formula}
Let $q\ge1$, $1\le p\le q$, and $\eta=p/q$. Choose a local $q$th
root $w=z^{1/q}$ and put $\omega=e^{2\pi i/q}$. Root-of-unity
filtering gives
\begin{equation}\label{dp:filter}
 \boxed{z^{p/q}\Phi(z,s,p/q)
       =q^{s-1}\sum_{\ell=0}^{q-1}
                     \omega^{-p\ell}\Li_s(\omega^\ell w).}
\end{equation}
To prove it, expand $\Li_s(\omega^\ell w)$ for $|w|<1$ and use
\[
 \frac1q\sum_{\ell=0}^{q-1}\omega^{\ell(k-p)}
 =\begin{cases}1,&k\equiv p\pmod q,\\0,&\text{otherwise}.
   \end{cases}
\]
The surviving positive indices are $k=qn+p$, including $p=q$ with
$n\ge0$, which yields \eqref{dp:filter}. The formula is entire in
the spectral order $s$ for $|z|<1$.

For $m\ge0$, write $\partial_s^m\Li_s$ for order derivatives, not
argument derivatives. Differentiation of \eqref{dp:filter} gives
\begin{align}
 z^{p/q}\partial_s^m\Phi(z,s,p/q)
 ={}&q^{s-1}\sum_{\ell=0}^{q-1}\omega^{-p\ell}
 \sum_{h=0}^m\binom mh(\log q)^{m-h}
              \partial_s^h\Li_s(\omega^\ell w).
 \label{dp:filterjets}
\end{align}
Thus every rational-shift subtraction term and every spectral jet
in \eqref{dp:primitiveeq} is a finite combination of classical
polylogarithms and their order derivatives. At integer resonances the
primitive orders are the even integers $2-2N+2j$; at the half-integer
points they are the odd integers $1-2M+2j$. Nonpositive integer
polylogarithms themselves reduce to rational functions by repeated
application of $z\,\mathrm d/\mathrm dz$ to $z/(1-z)$, while their
order derivatives retain additional information.

For a larger rational shift $\eta=L+p/q$, $L\ge0$ integer,
one must first apply the finite-head identity
\begin{equation}\label{dp:shift}
 z^\eta\Phi(z,s,\eta)
 =z^{p/q}\Phi(z,s,p/q)
       -\sum_{r=0}^{L-1}\frac{z^{r+p/q}}{(r+p/q)^s}.
\end{equation}
This follows by removing the first $L$ terms and is valid for every
$s$. Differentiating it retains the corresponding finite logarithmic
terms. Dropping this finite head would invalidate the claimed
polylogarithmic reduction for $\eta>1$.

For the quartic family $\eta=1/4$, the four colors are $1,i,-1,-i$:
\begin{equation}\label{dp:quarterfilter}
 z^{1/4}\Phi(z,s,1/4)=4^{s-1}
 \{\Li_s(w)-i\Li_s(iw)-\Li_s(-w)+i\Li_s(-iw)\},
 \quad w=z^{1/4}.
\end{equation}
Together with \eqref{dg:quarticshift} and
\eqref{dp:primitiveeq}, this gives a finite Gaussian-colored
polylogarithm description of the counterterms for every resonant
quartic moment and all its harmonic derivatives. It does not furnish
the missing exact relation certificates for the manuscript's
conjectural $S_6$ or $S_8$ evaluations.

\begin{remark}[Boundary values of the primitive]\label{dp:boundary}
In the strip of Theorem~\ref{dg:master}, the bracketed coefficients
in \eqref{dp:F} are absolutely summable. Dividing them by
$n+\eta$ preserves this property. Therefore the primitive itself has
a finite radial endpoint, obtainable from its one combined series.
At a resonant parameter, however, some separate terms on the right of
\eqref{dp:primitiveeq} can be singular at $z=1$. Their combined
expression, or its normally convergent coefficient series, is the
appropriate endpoint object. No separate assignment of values to
divergent polylogarithmic pieces is implicit.
\end{remark}

% END sections/03_polylog_primitives.tex

% BEGIN sections/04_nonlinear.tex
\section{Nonlinear coordinate changes of Stieltjes finite parts}
\label{nc:section}

The incoming Mellin--dilation and Stieltjes--harmonic reports explicitly
ask for the replacement of a local affine coordinate $qt$ by
$f(t)=qt+c_2t^2+\cdots$, including every Stieltjes index, every argument
derivative, and a composition law \cite{incoming-md,incoming-sh}.
The answer is a finite residue formula. Besides recovering their affine
polynomial, it yields a sharp additional phenomenon: if $f'(0)=1$, the
finite part of $\gamma_n^{(r)}$ transforms without a local correction
whenever $n\geq2r$.

The method below is direct analytic regularization. Its general
finite-part principle is not presented as a claim of historical priority.
The contribution to the stated source question is the complete formula,
its specialization to the Stieltjes hierarchy, the exact composition law,
and the sharp filtration of the surviving local terms.

\subsection{A universal residue formula}

All coordinates in this section are local real coordinates, with the
singular approach from the right. Let
\begin{equation}\label{nc:loggerm}
 h(x)=\sum_{\ell=0}^{L}a_\ell(x)(\log x)^\ell,\qquad x>0,
\end{equation}
where the $a_\ell$ are meromorphic germs at zero, each with a finite pole
order. Define $\FP_xh$ by the constant term of
$\int_\eps h(x)\phi(x)\dd x$ as $\eps\downarrow0$, for a smooth
compactly supported test function $\phi$. Every negative power of
$\eps$ and every nonconstant power of $\log\eps$ is discarded.
The positive-side germ can be extended by zero on the left; a fixed
smooth completion can equally well be added. Such a completion changes
neither the residue nor the coordinate correction.

Let $f$ be a real-analytic local diffeomorphism with
\begin{equation}\label{nc:coordinate}
 f(0)=0,\qquad f'(0)=q>0,
 \qquad \lambda_f(t)=\log\frac{f(t)}t.
\end{equation}
The last function is analytic at zero, with
$\lambda_f(0)=\log q$. Pullback here means scalar distributional
pullback:
\begin{equation}\label{nc:pullback}
 \langle f^*T,\phi\rangle
 =\left\langle T,
       (\phi\circ f^{-1})(f^{-1})'\right\rangle.
\end{equation}
In particular, the pullback of an ordinary function is its composition
with $f$.

\begin{theorem}[Nonlinear finite-part coordinate law]
\label{nc:universal}
For the germ \eqref{nc:loggerm}, define
\begin{equation}\label{nc:residue-anomaly}
 \left\langle\mathfrak A_f[h],\phi\right\rangle
 =\sum_{\ell=0}^{L}\frac1{\ell+1}
  \Res_{t=0}
    a_\ell(f(t))\lambda_f(t)^{\ell+1}\phi(t)\dd t.
\end{equation}
Then
\begin{equation}\label{nc:universal-law}
 \boxed{\quad
 f^*\FP_xh=\FP_t(h\circ f)+\mathfrak A_f[h].
 \quad}
\end{equation}
If every $a_\ell$ has pole order at most $R$, the correction is a
linear combination of
$\delta_0,\delta_0',\ldots,\delta_0^{(R-1)}$; for $R=0$ it is zero.
Every coefficient is determined by finite Taylor data. The residue
notation for a smooth $\phi$ means its finite Taylor coefficient;
holomorphic extension of the test function is unnecessary.
\end{theorem}

\begin{proof}
It suffices to take $h(x)=a(x)\log^\ell x$. For $\Re z$ sufficiently
large, the pairing of $x_+^za(x)$ with a test function is an ordinary
integral. If $a(x)=x^{-R}A(x)$ with $A$ holomorphic, subtract the Taylor
polynomial of $A(x)\phi(x)$ through degree $R-1$. The remainder is
locally integrable for $z$ near zero, uniformly on a small closed disk.
This constructs a meromorphic distribution with at most a simple pole
at zero, and
\begin{equation}\label{nc:regulator-fp}
 \FP_x\bigl(a(x)\log^\ell x\bigr)
       =\ell![z^\ell]\bigl(x_+^za(x)\bigr).
\end{equation}
Here the bracket is a Laurent coefficient. To check the normalization,
the only exceptional Taylor monomial is $x^{z-1}$, whose continued
integral to a fixed positive $h_0$ is $h_0^z/z$. Its coefficient
$[z^\ell]$ is $\log^{\ell+1}h_0/(\ell+1)!$, precisely the constant
remaining after subtracting the lower-endpoint logarithm. Every other
Taylor monomial has a denominator nonzero at $z=0$.

After pullback the regulated integrand is
\[
 t^z a(f(t))e^{z\lambda_f(t)}\phi(t).
\]
Its polar Taylor term at zero has numerator
\begin{equation}\label{nc:Bregulator}
 B(z)=\Res_{t=0}a(f(t))e^{z\lambda_f(t)}\phi(t)\dd t.
\end{equation}
Thus the exceptional contribution is $B(z)h_0^z/z$. On the other
hand, the cutoff integral of the same term is
$B(z)(h_0^z-\eps^z)/z$. In its coefficient $[z^\ell]$, the lower
endpoint has constant term $[z^{\ell+1}]B(z)$: all its other terms
are positive powers of $\log\eps$. Consequently the meromorphic
coefficient exceeds the coefficientwise cutoff constant by
\[
 \ell![z^{\ell+1}]B(z)
   =\frac1{\ell+1}\Res_{t=0}
      a(f(t))\lambda_f(t)^{\ell+1}\phi(t)\dd t.
\]
Nonpolar Taylor terms have lower-endpoint powers different from zero
and introduce no constant discrepancy. The subtracted remainder has
uniform integrable majorants, so Cauchy's coefficient formula justifies
the coefficient extraction there. This proves
\eqref{nc:universal-law}.

Finally, $a_\ell(f(t))$ has pole order at most $R$ and $\lambda_f$ is
analytic. Only derivatives $\phi^{(j)}(0)$ with $j\leq R-1$ can occur
in its residue. This proves both the support and finite-data claims.
\end{proof}

\subsection{Composition and the transformed singular germ}

\begin{theorem}[Exact residue cocycle]
\label{nc:composition}
For two orientation-preserving analytic coordinate germs $f,g$,
\begin{equation}\label{nc:cocycle}
 \boxed{\quad
 \mathfrak A_{g\circ f}[h]
    =f^*\mathfrak A_g[h]+\mathfrak A_f[h\circ g].
 \quad}
\end{equation}
The last term acts on the transformed germ $h\circ g$, which again has
the form \eqref{nc:loggerm}. This dependence is part of the composition
law.
\end{theorem}

\begin{proof}
The assertion follows by applying Theorem~\ref{nc:universal} twice.
There is also an explicit coefficient proof. For a single term
$h(x)=a(x)\log^\ell x$, put
$A=\lambda_f(t)$ and $B=\lambda_g(f(t))$. Then
\[
 \lambda_{g\circ f}=A+B,\qquad
 h(g(y))=\sum_{k=0}^\ell\binom\ell k
   a(g(y))\lambda_g(y)^{\ell-k}(\log y)^k.
\]
Residue substitution $y=f(t)$ in the pullback
\eqref{nc:pullback} shows that $f^*\mathfrak A_g[h]$ contributes the
factor $B^{\ell+1}/(\ell+1)$. The other term contributes
\[
 \sum_{k=0}^\ell\frac{\binom\ell k}{k+1}B^{\ell-k}A^{k+1}
 =\frac{(A+B)^{\ell+1}-B^{\ell+1}}{\ell+1}.
\]
The sum is precisely \eqref{nc:residue-anomaly} for $g\circ f$.
This finite polynomial identity proves composition with all scale and
nonlinear terms retained.
\end{proof}

For later implementation, the same residue convention gives
\begin{equation}\label{nc:delta-pullback}
 \langle f^*\delta_0^{(r)},\phi\rangle
   =(-1)^rr!\Res_{t=0}f(t)^{-r-1}\phi(t)\dd t.
\end{equation}
Indeed, substitution $x=f(t)$ changes the residue into the coefficient
of $x^r$ in $(\phi\circ f^{-1})(f^{-1})'$, in agreement with the
definition of $\delta_0^{(r)}$. Equations
\eqref{nc:cocycle} and \eqref{nc:delta-pullback} provide an exact
composition algorithm. A formula that retained $h$ unchanged in the
last term of \eqref{nc:cocycle} would generally be incorrect.

\subsection{Every Stieltjes index and argument derivative}

The local singularity is fixed by the Hurwitz convention:
\begin{equation}\label{nc:Stieltjes-local}
 \gamma_n(x)=\frac{\log^n x}{x}+\gamma_n(1+x),
 \qquad \gamma_0(x)=-\psi(x).
\end{equation}
The second term is analytic near zero. Set
\begin{equation}\label{nc:Tr}
 T_{n,r}=\FP_x\gamma_n^{(r)}(x),\qquad G_n=T_{n,0}.
\end{equation}
Locally these may be regarded as the periodic finite parts used in
\cite{incoming-md,incoming-sh}; the smooth completion in
\eqref{nc:Stieltjes-local} is then the same on both sides of the
endpoint. Every comparison involving $D^rG_n$ below uses this common
smooth periodic completion. In particular, the analytic remainder is
smooth across zero and contributes no step-function derivatives.

For $r\geq0$, let
\begin{align}
 \mathcal E_r(z)&=\prod_{j=1}^r(1-z/j),\qquad \mathcal E_0(z)=1,
       \label{nc:Er}\\
 b_{n,r}(L)&=n![z^{n+1}]e^{Lz}\mathcal E_r(z),\notag\\
 \alpha_{n,r}(L)&=b_{n,r}(L)-b_{n,r}(0)
       =n![z^{n+1}](e^{Lz}-1)\mathcal E_r(z).
       \label{nc:alpha}
\end{align}
If $e_j(r)$ is the elementary symmetric polynomial of degree $j$ in
$1,1/2,\ldots,1/r$, with $e_0(r)=1$ and $e_j(r)=0$ for $j>r$, then
\begin{equation}\label{nc:alpha-finite}
 \alpha_{n,r}(L)=n!\sum_{j=0}^{\min(r,n)}
      \frac{(-1)^je_j(r)}{(n+1-j)!}L^{n+1-j}.
\end{equation}
The subtraction of $b_{n,r}(0)$ is necessary because $T_{n,r}$ in
\eqref{nc:Tr} is the finite part of a pointwise derivative.

\begin{theorem}[Nonlinear Stieltjes transport]
\label{nc:Stieltjes}
For every $n,r\geq0$,
\begin{equation}\label{nc:Stieltjes-law}
 f^*T_{n,r}=\FP_t\gamma_n^{(r)}(f(t))+\mathfrak A_{n,r}(f),
\end{equation}
where
\begin{equation}\label{nc:Stieltjes-residue}
 \langle\mathfrak A_{n,r}(f),\phi\rangle
   =(-1)^rr!\Res_{t=0}
     f(t)^{-r-1}\alpha_{n,r}(\lambda_f(t))\phi(t)\dd t.
\end{equation}
In particular,
\begin{align}
 \mathfrak A_{n,r}(f)&=\sum_{j=0}^r c_{n,r,j}(f)\delta_0^{(j)},
       \label{nc:finite-delta}\\
 c_{n,r,j}(f)&=\frac{(-1)^{r+j}r!}{j!}
   [t^{r-j}]\left(\frac{f(t)}t\right)^{-r-1}
       \alpha_{n,r}\!\left(\log\frac{f(t)}t\right).
       \label{nc:coefficient}
\end{align}
The coordinate jet through degree $r+1$ suffices for every Stieltjes
index $n$.
\end{theorem}

\begin{proof}
Differentiate $x^{z-1}$ in its argument:
\[
 \partial_x^rx^{z-1}
    =(-1)^rr!\mathcal E_r(z)x^{z-r-1}.
\]
Taking $n![z^n]$ gives the exact singular polynomial
\begin{equation}\label{nc:singular-polynomial}
 \partial_x^r\frac{\log^n x}{x}
  =(-1)^rr!x^{-r-1}
      \sum_{j=0}^{\min(r,n)}
       \frac{(-1)^je_j(r)n!}{(n-j)!}\log^{n-j}x.
\end{equation}
Apply Theorem~\ref{nc:universal} term by term. Its factor
$1/(n-j+1)$ converts the displayed polynomial into
\eqref{nc:alpha-finite}, while the analytic remainder in
\eqref{nc:Stieltjes-local} contributes no residue. Expanding the test
function through degree $r$ yields \eqref{nc:coefficient}, since
$\langle\delta_0^{(j)},\phi\rangle=(-1)^j\phi^{(j)}(0)$.
The coefficient in that formula uses the expansion of $f(t)/t$ only
through degree $r$, proving the last assertion.
\end{proof}

For comparison, distributional differentiation gives
\begin{equation}\label{nc:derivative-contact}
 D^rG_n=T_{n,r}+b_{n,r}(0)\delta_0^{(r)},\qquad
 b_{n,r}(0)=(-1)^{n+1}n!e_{n+1}(r).
\end{equation}
This also follows by the regulator computation: at identity scale the
polar numerator for $\partial_x^rx^{z-1}$ is
$\mathcal E_r(z)\delta_0^{(r)}$. Thus
\begin{align}
 f^*D^rG_n&=\FP_t\gamma_n^{(r)}(f(t))+\mathfrak C_{n,r}(f),
       \label{nc:distributional-law}\\
 \langle\mathfrak C_{n,r}(f),\phi\rangle
   &=(-1)^rr!\Res_{t=0}
       f(t)^{-r-1}b_{n,r}(\lambda_f(t))\phi(t)\dd t.
       \label{nc:Cresidue}
\end{align}
The finite coefficient formula for $\mathfrak C$ is
\eqref{nc:coefficient} with $\alpha$ replaced by $b$.

\begin{corollary}[Recovery of the affine contact polynomial]
\label{nc:affine}
At $f(t)=qt$, $q>0$,
\begin{align}
 \mathfrak A_{n,r}(qt)&=q^{-r-1}\alpha_{n,r}(\log q)
                         \delta_0^{(r)},\label{nc:affine-A}\\
 D_t^r f^*G_n
   &=\FP_t\bigl(q^r\gamma_n^{(r)}(qt)\bigr)
        +\frac{b_{n,r}(\log q)}q\delta_0^{(r)}.
       \label{nc:affine-C}
\end{align}
The second identity is exactly the local form of the incoming
integer-dilation theorem, including its normalization by $1/q$.
\end{corollary}
\begin{proof}
The function $f(t)/t$ is constant. Only the $j=r$ term survives in
\eqref{nc:coefficient}; also
$D_t^rf^*G_n=q^rf^*D^rG_n$.
\end{proof}

\subsection{Explicit corrections through the second derivative}

Write
\begin{equation}\label{nc:f-expansion}
 f(t)=qt(1+a t+b t^2+O(t^3)),\qquad L=\log q.
\end{equation}
Thus $a=f''(0)/(2q)$ and $b=f'''(0)/(6q)$. The first coefficient
polynomials are
\begin{align}
 \alpha_{0,r}(L)&=L,\notag\\
 \alpha_{1,r}(L)&=\frac{L^2}{2}-H_rL,\label{nc:first-alpha}\\
 \alpha_{2,r}(L)&=\frac{L^3}{3}-H_rL^2
                    +(H_r^2-H_r^{(2)})L.\notag
\end{align}
Here $H_0=H_0^{(2)}=0$. At $r=0$, the entire answer is
\begin{equation}\label{nc:undifferentiated}
 \mathfrak A_{n,0}(f)=\frac{L^{n+1}}{q(n+1)}\delta_0.
\end{equation}
Nonlinear coefficients first affect argument derivatives.

For $r=1$ set $\alpha=\alpha_{n,1}(L)$; primes below differentiate
this polynomial in $L$. Equation~\eqref{nc:coefficient} becomes
\begin{equation}\label{nc:r1}
 \mathfrak A_{n,1}(f)
  =q^{-2}\left[\alpha\delta_0'
                   +a(2\alpha-\alpha')\delta_0\right].
\end{equation}
For $r=2$, with $\alpha=\alpha_{n,2}(L)$, it gives
\begin{align}\label{nc:r2}
 \mathfrak A_{n,2}(f)=q^{-3}\bigl[&\alpha\delta_0''
       +2a(3\alpha-\alpha')\delta_0'\notag\\
       &+\{(12a^2-6b)\alpha+(2b-7a^2)\alpha'
                                      +a^2\alpha''\}\delta_0\bigr].
\end{align}
To verify these formulas directly, use
\[
 \lambda_f=L+a t+(b-a^2/2)t^2+O(t^3),\qquad
 (1+a t+b t^2)^{-3}
      =1-3a t+(6a^2-3b)t^2+O(t^3).
\]
The same displays with $b_{n,r}$ in place of $\alpha_{n,r}$ give the
correction $\mathfrak C_{n,r}$ for distributional derivatives.

At unit tangent, these are particularly short:
\begin{equation}\label{nc:r1unit}
 \mathfrak A_{n,1}=
 \begin{cases}
  -a\delta_0,&n=0,\\
  a\delta_0,&n=1,\\
  0,&n\geq2,
 \end{cases}
\end{equation}
and
\begin{equation}\label{nc:r2unit}
 \mathfrak A_{n,2}=
 \begin{cases}
  -2a\delta_0'+(2b-7a^2)\delta_0,&n=0,\\
  3a\delta_0'+(-3b+\tfrac{23}{2}a^2)\delta_0,&n=1,\\
  -2a\delta_0'+(2b-10a^2)\delta_0,&n=2,\\
  3a^2\delta_0,&n=3,\\
  0,&n\geq4.
 \end{cases}
\end{equation}
These expressions retain curvature and the cubic coordinate coefficient
even though the linear scale has disappeared.

\begin{example}[A direct cutoff check of the curvature terms]
\label{nc:direct-cutoff}
Take $f(t)=t+a t^2+b t^3+O(t^4)$ and $\gamma_0''(x)$, whose singular
part is $2x^{-3}$. Its inverse is
$g(x)=x-a x^2+(2a^2-b)x^3+O(x^4)$. For a test function equal to one
near zero, pullback uses the inverse Jacobian
$g'(x)=1-2a x+(6a^2-3b)x^2+O(x^3)$.
A primitive of the terms that can contribute a finite constant is
\[
 -x^{-2}+4a x^{-1}+(12a^2-6b)\log x.
\]
It has zero constant in the $x$ coordinate. After $x=f(\eps)$ its
constant is $(2b-3a^2)-4a^2=2b-7a^2$, giving the $\delta_0$
coefficient in \eqref{nc:r2unit}. For a test function equal to $t$,
the inverse multiplier is $g(x)g'(x)=x-3a x^2+O(x^3)$; the relevant
primitive is $-2/x-6a\log x$. Its composed constant is $2a$, which
equals the action of $-2a\delta_0'$ on that test function.
This calculation uses cutoff primitives directly and confirms both
nonlinear signs without invoking the residue formula.
\end{example}

\subsection{Sharp invariance and the order of the local correction}

\begin{theorem}[Sharp Stieltjes-index filtration]
\label{nc:rigidity}
Suppose $f'(0)=1$, put $a=f''(0)/2$, and let $r\geq1$. Then
$\mathfrak A_{n,r}(f)$ is a linear combination of
$\delta_0^{(j)}$ with
\begin{equation}\label{nc:order-filtration}
 0\leq j\leq\min\{r-1,\,2r-n-1\}.
\end{equation}
An empty range means the zero distribution. In particular,
\begin{equation}\label{nc:high-invariance}
 \boxed{\quad
 f^*\FP_x\gamma_n^{(r)}
      =\FP_t\gamma_n^{(r)}(f(t))\qquad(n\geq2r).
 \quad}
\end{equation}
For $r\leq n\leq2r-1$, let
$d=n+1-r$ and $k=2r-n-1$. The coefficient of the highest possible
derivative is
\begin{equation}\label{nc:sharp-coefficient}
 [\delta_0^{(k)}]\mathfrak A_{n,r}(f)
        =\frac{(-1)^kn!}{d!\,k!}a^d.
\end{equation}
For $0\leq n\leq r$, the coefficient of $\delta_0^{(r-1)}$ is
\begin{equation}\label{nc:lower-sharp}
 [\delta_0^{(r-1)}]\mathfrak A_{n,r}(f)
       =(-1)^{n+1}r\,n!e_n(r)a.
\end{equation}
The two formulas agree at $n=r$. If $a\ne0$, every stated upper
order is attained. At the last nonzero index,
\begin{equation}\label{nc:last-index}
 \mathfrak A_{2r-1,r}(f)
     =\frac{(2r-1)!}{r!}a^r\delta_0.
\end{equation}
For $r=0$, every correction vanishes at unit tangent by
\eqref{nc:undifferentiated}.
\end{theorem}

\begin{proof}
Now $\lambda_f(t)=a t+O(t^2)$. Formula~\eqref{nc:alpha-finite}
shows that the smallest possible power of $L$ in $\alpha_{n,r}$ is
$\max(1,n+1-r)$. Therefore the analytic factor multiplying
$t^{-r-1}$ in \eqref{nc:Stieltjes-residue} vanishes to at least that
order. Its residue against $\phi$ can use only derivatives of orders
at most $r-\max(1,n+1-r)$, proving
\eqref{nc:order-filtration} and \eqref{nc:high-invariance}.

For $r\leq n\leq2r-1$ the lowest term is, using $e_r(r)=1/r!$,
\[
 \alpha_{n,r}(L)
    =\frac{(-1)^rn!}{r!\,d!}L^d+O(L^{d+1}).
\]
To produce $\delta_0^{(k)}$ in \eqref{nc:coefficient}, only the
coefficient of $t^{r-k}=t^d$ is required. It is the displayed
coefficient times $a^d$; the constant term of $(f(t)/t)^{-r-1}$ is
one. Substitution proves \eqref{nc:sharp-coefficient}. For $n\leq r$
the linear coefficient of $\alpha_{n,r}$ is
$(-1)^nn!e_n(r)$. Taking $j=r-1$ in the same formula proves
\eqref{nc:lower-sharp}. All these elementary symmetric coefficients
are nonzero in their indicated ranges. Equation~\eqref{nc:last-index}
is the case $d=r,k=0$.
\end{proof}

\begin{remark}[Higher tangency gives stronger invariance]
\label{nc:higher-tangency}
If $f(t)=t+O(t^{d+1})$, then $\lambda_f=O(t^d)$. If $d>r$, every
$\mathfrak A_{n,r}$ vanishes. If $1\leq d\leq r$, it suffices that
\[
 n\geq r+\lfloor r/d\rfloor.
\]
This threshold is sharp in general. For
$f(t)=t+c t^{d+1}+O(t^{d+2})$, set
$h=\lfloor r/d\rfloor$, $n=r+h-1$ and $k=r-dh$. The coefficient of
$\delta_0^{(k)}$ is
$(-1)^kn!c^h/(h!k!)$. This follows from the same lowest-power argument.
\end{remark}

\subsection{Polygamma functions, full derivatives, and pointwise products}

For polygamma, the correction in
\eqref{nc:Stieltjes-law} is $-\mathfrak A_{0,r}$ because
$\gamma_0^{(r)}=-\psi^{(r)}$. For example, at
$f(t)=t+a t^2+O(t^3)$,
\begin{equation}\label{nc:trigamma}
 f^*\FP_x\psi'(x)
    =\FP_t\psi'(f(t))+a\delta_0.
\end{equation}
The undifferentiated function $\log\Gamma(x)$ is locally integrable;
its logarithmic singularity has no negative-power coefficient in
\eqref{nc:loggerm}. Its ordinary pullback therefore has no anomaly.

A derivative in the argument of $\gamma_n$ differs from a full
derivative of $\gamma_n(f(t))$. The latter is also finite and explicit.
Let $B_{k,r}$ denote the partial exponential Bell polynomial in the
ordinary Fa\`a di Bruno formula. For $k\geq1$,
\begin{align}\label{nc:chain-rule}
 D_t^kf^*G_n-&\FP_tD_t^k\bigl(\gamma_n(f(t))\bigr)\notag\\
   &=\sum_{r=1}^k
       B_{k,r}(f',f'',\ldots,f^{(k-r+1)})\,
                         \mathfrak C_{n,r}(f).
\end{align}
Indeed, the distributional chain rule holds first for smooth functions
and then by pullback and differentiation of distributions; it gives
the displayed Bell sum with $f^*D^rG_n$. Finite parts commute with
multiplication by a smooth function, directly from their test-function
definition, so \eqref{nc:distributional-law} proves
\eqref{nc:chain-rule}. To put the answer in constant delta coordinates,
use the finite identity
\begin{equation}\label{nc:delta-multiplier}
 v\delta_0^{(j)}
   =\sum_{\ell=0}^j(-1)^{j-\ell}\binom j\ell
       v^{(j-\ell)}(0)\delta_0^{(\ell)}.
\end{equation}
Thus no indefinite distribution operation remains in the formula.

The universal theorem also treats a pointwise product once its own
finite part has been specified. This gives a small exact illustration
relevant to collision normalizations.

\begin{corollary}[Coordinate law for a coincident Stieltjes product]
\label{nc:coincident}
Let $h(x)=\gamma_m(x)\gamma_n(x)$, $M=m+n$, and use
\eqref{nc:f-expansion}. Then
\begin{align}\label{nc:product-anomaly}
 \mathfrak A_f[h]={}&-\frac{L^{M+1}}{q^2(M+1)}\delta_0'\notag\\
 &+\biggl[
  \frac{a}{q^2}\left(L^M-\frac{2L^{M+1}}{M+1}\right)
  +\frac{\gamma_n(1)L^{m+1}}{q(m+1)}
  +\frac{\gamma_m(1)L^{n+1}}{q(n+1)}
                  \biggr]\delta_0.
\end{align}
As usual $L^0=1$, also when $L=0$. At unit tangent the sole nonzero
case is $m=n=0$, with correction $a\delta_0$.
\end{corollary}
\begin{proof}
Equation~\eqref{nc:Stieltjes-local} gives the complete nonintegrable
part
\[
 \frac{\log^{M}x}{x^2}
 +\frac{\gamma_n(1+x)\log^m x+
         \gamma_m(1+x)\log^n x}{x}.
\]
For the first term, expand $f^{-2}\lambda_f^{M+1}/(M+1)$ through
the coefficient of $t^{-1}$. It gives the first line of
\eqref{nc:product-anomaly} and its $a/q^2$ term. The two simple-pole
terms require only the values $\gamma_n(1)$ and $\gamma_m(1)$,
and yield the remaining terms. The analytic product remainder has no
residue.
\end{proof}

This corollary compares coordinate choices for a specified finite part
of an ordinary pointwise product. It does not identify that finite
part with a collision limit of separated products, and it does not
define multiplication of arbitrary distributions at common singular
supports. It supplies explicit local normalization data for such a
future comparison.

\subsection{Exact algorithm and independent checks}

For fixed $n,r$, the complete computation consists of four finite
operations: form $\mathcal E_r$, extract $\alpha_{n,r}$, truncate
$\log(f(t)/t)$ and $(f(t)/t)^{-r-1}$ through degree $r$, and read the
coefficients in \eqref{nc:coefficient}. The coefficients are rational
polynomials in the normalized coordinate derivatives and $\log q$,
with the overall factor $q^{-r-1}$. Composition is checked by the
binomial identity in Theorem~\ref{nc:composition}.

The accompanying program \src{code/verify_nonlinear.py} performs exact
rational symbolic checks of the coefficient generator, the scale-aware
cocycle, the generic first- and second-derivative formulas, and the sharp
filtration. It also compares the correction with an independent direct
cutoff-primitive calculation for finite test jets. The record is
\src{results/nonlinear_checks.json}. These finite tests check the
implementation and normalization; the all-index analytic assertions
are established by the proofs above.

% END sections/04_nonlinear.tex

% BEGIN sections/05_transport.tex
\section{Three unequal frequencies with exact endpoint normalization}
\label{tr:section}

The incoming three-factor formula treats translated copies of the same
periodic Hurwitz function. Its next question allows three distinct positive
integer frequencies \cite{incoming-sh}. Solving the Fourier constraint
directly produces weighted linear forms. A finite multiplication identity
gives a shorter route: first replace each integer-frequency factor by
finitely many unit-frequency translates. All local scales can then be
accounted for before the product is formed.

\subsection{Separated singular grids and a finite lifting identity}
Fix $p_1,p_2,p_3\in\N_{>0}$ and $0\leq a_j<1$. Put
\begin{equation}\label{tr:grids}
 \mathcal B_j=\left\{b_{j,h}=\frac{a_j+h}{p_j}:
                       0\leq h<p_j\right\}\subset[0,1).
\end{equation}
Assume the three sets $\mathcal B_j$ are pairwise disjoint. The singular
points of the $j$th original factor are $x=-b$ modulo one, with
$b\in\mathcal B_j$; hence this assumption is precisely disjointness of
the singular grids. It is directly decidable for rational data. In a
pairwise arithmetic form, if $d=\gcd(p_i,p_j)$, it says
\[
 (p_i/d)a_j-(p_j/d)a_i\notin\Z\qquad(i\ne j).
\]
The grid formulation will make the local geometry transparent.

\begin{lemma}[Finite periodic multiplication]\label{tr:multiplication}
For a positive integer $p$, $0\leq a<1$, and an argument away from the
singular grid,
\begin{equation}\label{tr:zeta-lift}
 \zeta(s,\{px+a\})
  =p^{-s}\sum_{h=0}^{p-1}
      \zeta\!\left(s,\left\{x+\frac{a+h}{p}\right\}\right).
\end{equation}
Both sides agree as meromorphic functions of $s$. If $L=\log p$, then
for every $m\geq0$,
\begin{align}\label{tr:gamma-lift}
 \gamma_m(\{px+a\})={}&\frac1p\sum_{h=0}^{p-1}
      \sum_{\ell=0}^m\binom m\ell L^{m-\ell}
      \gamma_\ell\!\left(\left\{x+\frac{a+h}{p}\right\}\right)
       -\frac{L^{m+1}}{m+1}.
\end{align}
\end{lemma}
\begin{proof}
Write $y=\{px+a\}$. The $p$ numbers
$\{x+(a+h)/p\}$, counted with multiplicity, are a permutation of
$(y+k)/p$, $0\leq k<p$. Splitting the defining Hurwitz series into
residue classes gives
$\sum_k\zeta(s,(y+k)/p)=p^s\zeta(s,y)$ for $\Re s>1$.
Meromorphic continuation proves \eqref{tr:zeta-lift}.

Set $s=1-z$. The expansion
$\zeta(1-z,y)=-1/z+\sum_{m\geq0}\gamma_m(y)z^m/m!$
shows that the nonpolar part on the right of \eqref{tr:zeta-lift} is
\[
 \frac{e^{Lz}}p\sum_{h=0}^{p-1}
       \sum_{\ell\geq0}\gamma_\ell(\{x+(a+h)/p\})\frac{z^\ell}{\ell!}
       -\frac{e^{Lz}-1}{z}.
\]
Its $m$th derivative at zero is \eqref{tr:gamma-lift}. This derivation
also fixes the sign and the absence of a factor $1/p$ in the last term.
\end{proof}

Let $G_m$ be the unit-coordinate periodic finite part of $\gamma_m$.
It has mean zero. For completeness,
\[
 \gamma_m(x)=x^{-1}\log^m x+\gamma_m(1+x)
\]
has a singular term of finite-part integral zero on $(0,1)$. Also
$\int_1^2\zeta(1-z,x)\dd x=-1/z$, by
$\partial_x\zeta(-z,x)=z\zeta(1-z,x)$ and the Hurwitz shift relation.
Its regular coefficients give $\int_1^2\gamma_m(x)\dd x=0$.
Therefore $\langle G_m,1\rangle=0$.

Taking the constant term of the pointwise identity
\eqref{tr:gamma-lift} in the same ambient coordinate at every singular
point gives an equality of distributions. Define
\begin{align}
 c_m(p)&=\frac{\log^{m+1}p}{m+1},\notag\\
 W_{p,a;m}&=\frac1p\sum_{h=0}^{p-1}\sum_{\ell=0}^m
       \binom m\ell(\log p)^{m-\ell}\,
              G_\ell(x+(a+h)/p).
       \label{tr:W}
\end{align}
Then
\begin{equation}\label{tr:ambient-factor}
 \FP_x\gamma_m(\{px+a\})=W_{p,a;m}-c_m(p).
\end{equation}
In particular,
\begin{equation}\label{tr:mean}
 \FP_x\int_0^1\gamma_m(\{px+a\})\dd x
       =-\frac{\log^{m+1}p}{m+1}.
\end{equation}
This nonzero mean is a simple check on the endpoint scale. It would be
missed by replacing a fixed-coordinate finite part by an uncorrected
pullback of a mean-zero periodic distribution.

\subsection{A complete all-index transport formula}
For pairwise distinct unit-frequency shifts $b_1,b_2,b_3$, let
\begin{equation}\label{tr:Idef}
 I_{\boldsymbol\ell}(\boldsymbol b)
   =\FP_x\int_0^1\prod_{j=1}^3
              \gamma_{\ell_j}(\{x+b_j\})\dd x.
\end{equation}
For $0<c<1$, write
\begin{equation}\label{tr:Jdef}
 J_{uv}(c)=\FP_x\int_0^1
           \gamma_u(x)\gamma_v(\{x+c\})\dd x.
\end{equation}
All finite parts use the right coordinate $x-x_0$ at each singularity.
Different cutoffs may tend independently to zero. Since the singular
supports are disjoint, their distribution products are defined by
ordinary multiplication by smooth functions near each singularity;
their integrals equal these finite parts.

Fix $m_1,m_2,m_3\geq0$, and abbreviate
\[
 L_j=\log p_j,\quad c_j=\frac{L_j^{m_j+1}}{m_j+1},\quad
 A_{j,\ell}=\frac1{p_j}\binom{m_j}{\ell}L_j^{m_j-\ell}.
\]
As usual, a zeroth power is one, also when $L_j=0$.

\begin{theorem}[All-index unequal-frequency reduction]
\label{tr:finite-transport}
Under the disjoint-grid hypothesis, the ambient-coordinate finite part
\[
 I^{\boldsymbol p}_{\boldsymbol m}(\boldsymbol a)
 =\FP_x\int_0^1\prod_{j=1}^3
                 \gamma_{m_j}(\{p_jx+a_j\})\dd x
\]
has the finite expression
\begin{align}
 I^{\boldsymbol p}_{\boldsymbol m}(\boldsymbol a)
 ={}&\sum_{h_1=0}^{p_1-1}\sum_{h_2=0}^{p_2-1}\sum_{h_3=0}^{p_3-1}
       \sum_{\ell_1=0}^{m_1}\sum_{\ell_2=0}^{m_2}\sum_{\ell_3=0}^{m_3}
       \left(\prod_{j=1}^3A_{j,\ell_j}\right)
       I_{\boldsymbol\ell}(b_{1,h_1},b_{2,h_2},b_{3,h_3})
       \notag\\
 &-\sum_{k=1}^3c_k
    \sum_{h_i=0}^{p_i-1}\sum_{h_j=0}^{p_j-1}
    \sum_{u=0}^{m_i}\sum_{v=0}^{m_j}
        A_{i,u}A_{j,v}
        J_{uv}(\{b_{j,h_j}-b_{i,h_i}\})
       -c_1c_2c_3 .\label{tr:finite-formula}
\end{align}
In the $k$th term of the second line, $i<j$ are the other two indices.
Every pair argument lies in $(0,1)$, and every triple has distinct
arguments. Thus no collision continuation is hidden in this formula.
\end{theorem}
\begin{proof}
Apply \eqref{tr:ambient-factor} to each of the three factors. Since
different grids have disjoint singular support, their product can be
expanded distributively, with no undefined distribution product. The
term containing three $W$ factors gives the first line of
\eqref{tr:finite-formula}. The three terms containing two $W$ factors
give its second line. The terms containing one $W$ factor integrate to
zero because each $G_\ell$ has zero mean. The constant term is
$-c_1c_2c_3$. This proves the formula. Equivalently, a partition of unity
at the separated endpoints proves the same identity directly by linearity
of the cutoff constant terms.
\end{proof}

This theorem resolves the stated unequal-frequency question at every
Stieltjes index. It permits a redundant finite list of coordinates rather
than claiming a minimal period basis. Integer frequencies are essential
to the finite multiplication step. Colliding grids require a separately
specified pointwise product finite part or joint regularization.

\subsection{The bilinear terms are finite Stieltjes expressions}
The pair terms in \eqref{tr:finite-formula} are already reduced in the
incoming reports \cite{incoming-sh,incoming-md}. We record their complete
coefficient formula so that the triple transport is constructive without
an unspecified lower-order operation. Put
\begin{align}
 \mathcal G(z;c)&=\zeta(1-z,c)+1/z,
 & U(z,w)&=\frac{\Gamma(1+z)\Gamma(1-z-w)}{\Gamma(1-w)},\notag\\
 V(z,w)&=\frac{U(z,w)-1}{z(z+w)},
 &r&=z+w.\label{tr:pair-generators}
\end{align}
Both $\mathcal G$ and $V$ are holomorphic at the indicated origin. For
$V$, the numerator vanishes on $z=0$ and on $z=-w$; these are distinct
linear factors. The formula is
\begin{align}
 \sum_{u,v\geq0}J_{uv}(c)\frac{z^uw^v}{u!v!}
 ={}&\frac{\mathcal G(r;c)-\mathcal G(w;c)}z
     +\frac{\mathcal G(r;1-c)-\mathcal G(z;1-c)}w\notag\\
 &-V(z,w)-V(w,z)\notag\\
 &+r\{V(z,w)\mathcal G(r;c)
           +V(w,z)\mathcal G(r;1-c)\}.\label{tr:pair-closure}
\end{align}
Every apparent singularity is removable. In particular,
\begin{equation}\label{tr:J00}
 J_{00}(c)=\gamma_1(c)+\gamma_1(1-c)-2\zeta(2).
\end{equation}

Here is a short derivation. Fourier expansion and Gamma reflection give
the meromorphic identity
\begin{align*}
 \int_0^1\zeta(1-z,x)\zeta(1-w,\{x+c\})\dd x
 ={}&B(z,1-z-w)\zeta(1-z-w,c)\\
 &+B(w,1-z-w)\zeta(1-z-w,1-c).
\end{align*}
It first holds in an absolutely convergent Fourier domain and extends
to $\Re z,\Re w>0$ by endpoint integrability. Its Laurent decomposition
at zero is
\[
 -\frac1{zw}+\frac{\mathcal G(w;c)}z
      +\frac{\mathcal G(z;1-c)}w
      +\sum_{u,v\geq0}J_{uv}(c)\frac{z^uw^v}{u!v!}.
\]
One obtains this decomposition by subtracting the singular Taylor term
at each of the two disjoint endpoints. Substitute
$B(z,1-r)=U(z,w)/z$ and $U=1+zrV$ in the beta expression and subtract
the three displayed polar terms. The remainder is exactly
\eqref{tr:pair-closure}.
Finally,
\[
 U(z,w)=\exp\!\left(\sum_{k\geq2}\frac{\zeta(k)}k
             \{(-1)^kz^k+(z+w)^k-w^k\}\right)
\]
shows that any requested pair coefficient uses only finitely many
Stieltjes constants and polynomial combinations of positive integer
zeta values. No infinite coefficient search is required.

\subsection{Ordinary colored Tornheim kernels suffice}
For $|X|=|Y|=1$, $X,Y\ne1$, define initially in an absolutely convergent
region
\begin{equation}\label{tr:Tdef}
 \mathcal T(A,B,C;X,Y)
     =\sum_{m,n\geq1}\frac{X^mY^n}{m^An^B(m+n)^C}.
\end{equation}
Its derivatives in $A,B,C$ are called spectral derivatives. They are
not derivatives of the fixed colors.

\begin{lemma}[Entire continuation and exact negative-order values]
\label{tr:T-entire}
The kernel \eqref{tr:Tdef} has a joint entire continuation in
$(A,B,C)$. For $\Re C>0$ it is given by
\begin{equation}\label{tr:T-mellin}
 \mathcal T(A,B,C;X,Y)=\frac1{\Gamma(C)}
   \int_0^\infty t^{C-1}\Li_A(Xe^{-t})\Li_B(Ye^{-t})\dd t.
\end{equation}
For every $N\geq0$,
\begin{equation}\label{tr:T-negative}
 \mathcal T(A,B,-N;X,Y)
   =\sum_{j=0}^N\binom Nj\Li_{A-j}(X)\Li_{B-N+j}(Y).
\end{equation}
\end{lemma}
\begin{proof}
The Gamma integral for $(m+n)^{-C}$ proves \eqref{tr:T-mellin} in a
common absolutely convergent domain. Since $X,Y\ne1$, each
polylogarithm is analytic in $t$ in a neighborhood of zero and entire
in its order. Their product has a Taylor expansion
$f(t)=\sum_{j=0}^Mf_jt^j+O_K(t^{M+1})$, uniformly on compact order
sets. Split the integral at one, subtract this polynomial on $(0,1)$,
and add $\sum_{j=0}^Mf_j/(C+j)$. The resulting expression extends to
$\Re C>-M-1$ after multiplication by $1/\Gamma(C)$. The simple Gamma
zeros cancel the apparent poles $C=0,-1,\ldots,-M$. The integral at
infinity is locally normal because of exponential decay. Increasing
$M$ proves joint entire continuation. At $C=-N$ the value is
$(-1)^Nf^{(N)}(0)$; applying
$\partial_t\Li_A(Xe^{-t})=-\Li_{A-1}(Xe^{-t})$ and Leibniz's rule
gives \eqref{tr:T-negative}.
\end{proof}

For pairwise distinct $b_1,b_2,b_3$, let $i<j$ be the indices other
than $k$ and put
\[
 X_k=e(b_i-b_k),\quad Y_k=e(b_j-b_k),\quad
 \theta_k=\frac\pi2(\lambda_i+\lambda_j-\lambda_k),
 \qquad e(t)=e^{2\pi it}.
\]
Define the finite sum
\begin{align}\label{tr:Sdef}
 S(\boldsymbol\lambda;\boldsymbol b)
 =\sum_{k=1}^3\{&e^{-i\theta_k}
      \mathcal T(\lambda_i,\lambda_j,\lambda_k;X_k,Y_k)\notag\\
 &+e^{i\theta_k}
      \mathcal T(\lambda_i,\lambda_j,\lambda_k;X_k^{-1},Y_k^{-1})\}.
\end{align}
The signs specify all six regions of the three-frequency hyperplane.
The identity rederived next is the unit-frequency result already
established in the incoming Stieltjes--harmonic report.

\begin{proposition}[The unit-frequency spectral identity]
\label{tr:unit-spectral}
For $\Re\lambda_j>0$,
\begin{equation}\label{tr:unit-K}
 K_1(\boldsymbol\lambda;\boldsymbol b)
 :=\int_0^1\prod_{j=1}^3\zeta(1-\lambda_j,\{x+b_j\})\dd x
  =\frac{\prod_j\Gamma(\lambda_j)}{(2\pi)^{\sum_j\lambda_j}}
       S(\boldsymbol\lambda;\boldsymbol b).
\end{equation}
Consequently
\begin{align}\label{tr:triple-coeff}
 I_{\boldsymbol\ell}(\boldsymbol b)
 =\left(\prod_{j=1}^3\ell_j!\right)
 [\lambda_1^{\ell_1+1}\lambda_2^{\ell_2+1}\lambda_3^{\ell_3+1}]
 \left\{\prod_{j=1}^3\frac{\Gamma(1+\lambda_j)}{(2\pi)^{\lambda_j}}
 S(\boldsymbol\lambda;\boldsymbol b)\right\}.
\end{align}
\end{proposition}
\begin{proof}
For $\Re\lambda_j>1$, the Fourier coefficients of each factor are
\[
 \widehat Z_{1-\lambda}(n)
  =\frac{\Gamma(\lambda)}{(2\pi|n|)^\lambda}
       e^{-i\pi\lambda\sgn(n)/2},\quad n\ne0,
 \qquad \widehat Z_{1-\lambda}(0)=0.
\]
This is Hurwitz's classical Fourier formula
\cite[Eq.~25.11.9]{dlmf-hurwitz}. Absolute convergence permits
termwise integration. The surviving triples satisfy
$n_1+n_2+n_3=0$, with every $n_j\ne0$. Either the unique negative
index is $k$, giving $n_i=m,n_j=n,n_k=-m-n$, or all signs are
reversed. Their denominators and phases are precisely
\eqref{tr:Sdef}. Near each separated singular point, only one factor
behaves like $t^{\lambda_j-1}$; the others are smooth. Local integrable
majorants extend the identity to $\Re\lambda_j>0$.

To extract the finite parts, split at disjoint neighborhoods of the
three singularities. The singular factor is
$t^{\lambda_k-1}+\zeta(1-\lambda_k,1+t)$. For a smooth multiplier
$g$, the continued integral is
\[
 \frac{g(0)h^{\lambda_k}}{\lambda_k}
 +\int_0^h t^{\lambda_k-1}(g(t)-g(0))\dd t.
\]
The second term is holomorphic near zero after clearing the at most
simple poles in the other order variables. The Laurent coefficient
$[\lambda_k^{\ell_k}]$ of the first term is the unit-coordinate
finite part divided by $\ell_k!$. Taylor subtraction gives uniform
remainders $O(\eps^{1-\rho})$ on a closed polydisk
$|\lambda_j|\leq\rho<1$ after multiplication by
$\lambda_1\lambda_2\lambda_3$. Cauchy's formula therefore justifies
simultaneous coefficient extraction. Finally,
$\lambda_j\Gamma(\lambda_j)=\Gamma(1+\lambda_j)$ converts
\eqref{tr:unit-K} to \eqref{tr:triple-coeff}.
\end{proof}

Combining Theorem~\ref{tr:finite-transport},
\eqref{tr:pair-closure}, and \eqref{tr:triple-coeff} is the promised
finite formula in standard colored kernels. Its largest triple spectral
derivative order is at most $m_1+m_2+m_3+3$.

\begin{corollary}[Spectral transport before taking finite parts]
\label{tr:spectral-lift}
For $\Re\lambda_j>0$,
\begin{align}\label{tr:Kp}
 K_{\boldsymbol p}(\boldsymbol\lambda;\boldsymbol a)
 &:=\int_0^1\prod_{j=1}^3
           \zeta(1-\lambda_j,\{p_jx+a_j\})\dd x\notag\\
 &=\left(\prod_{j=1}^3p_j^{\lambda_j-1}\right)
   \sum_{h_1=0}^{p_1-1}\sum_{h_2=0}^{p_2-1}\sum_{h_3=0}^{p_3-1}
       K_1(\boldsymbol\lambda;b_{1,h_1},b_{2,h_2},b_{3,h_3}).
\end{align}
Both sides have the same meromorphic continuation. The spectral
coefficient at the Stieltjes pole must still be converted to the named
ambient-coordinate finite part; \eqref{tr:finite-formula} performs
that conversion explicitly.
\end{corollary}
\begin{proof}
Multiply three instances of \eqref{tr:zeta-lift} with $s=1-\lambda_j$,
then integrate. The sums are finite and every resulting triple is
separated. Proposition~\ref{tr:unit-spectral} provides the continuation.
\end{proof}

% END sections/05_transport.tex

% BEGIN sections/06_correlations.tex
\section{Explicit polylogarithmic and Gamma correlation identities}
\label{tc:section}

The preceding transport theorem is an all-index result. This section
makes several consequences explicit, including an ordinary convergent
Gamma integral and a root-of-unity example whose lower terms reduce to
classical Stieltjes constants and logarithms.

\subsection{Finite arithmetic reduction of weighted Tornheim kernels}
The weighted double kernel that appears in a direct Fourier treatment
is itself a finite transform of \eqref{tr:Tdef}. Let $p,q\geq1$ be
integers and define
\begin{equation}\label{tc:weighted}
 \mathcal T_{p,q}(A,B,C;X,Y)
   =\sum_{m,n\geq1}\frac{X^mY^n}{m^An^B(pm+qn)^C}
\end{equation}
in an absolutely convergent region. Fix any roots $\xi^p=X$ and
$\eta^q=Y$.

\begin{proposition}[A finite root filter]\label{tc:root-filter}
If $|X|=|Y|=1$ and $X,Y\ne1$, then
\begin{align}\label{tc:weighted-reduction}
 \mathcal T_{p,q}(A,B,C;X,Y)
 =p^{A-1}q^{B-1}
   \sum_{h=0}^{p-1}\sum_{k=0}^{q-1}
     \mathcal T(A,B,C;\xi e(h/p),\eta e(k/q)).
\end{align}
The expression is independent of the chosen roots and is entire in
$(A,B,C)$. In particular, for $N\geq0$,
\begin{align}\label{tc:weighted-negative}
 \mathcal T_{p,q}(A,B,-N;X,Y)
 =\sum_{j=0}^N\binom Nj p^jq^{N-j}
             \Li_{A-j}(X)\Li_{B-N+j}(Y).
\end{align}
All $A$- and $B$-derivatives of this last identity are valid.
\end{proposition}
\begin{proof}
Substitute $M=pm$ and $N'=qn$ in \eqref{tc:weighted}. The indicator
of $p\mid M$ is $p^{-1}\sum_he(hM/p)$, and similarly for $q\mid N'$.
The factors $m^{-A}=p^AM^{-A}$ and $n^{-B}=q^B(N')^{-B}$ give
\eqref{tc:weighted-reduction}. Changing a chosen root permutes the
finite sum. No new color equals one because $X,Y\ne1$; hence
Lemma~\ref{tr:T-entire} proves entire continuation. Expanding
$(pm+qn)^N$ in the original convergent region proves
\eqref{tc:weighted-negative}, and continuation gives the full claim.
\end{proof}

Finite congruence conditions on the two indices can be imposed by a
further finite root filter. This explains why integral weights and
congruence restrictions arising from the unequal-frequency lattice can
be handled without introducing a new family of transcendental kernels.
It does not assert that the finite list produced in this way has no
further relations.

The derivative in the \emph{third} order requires separate data.
Differentiating an identity known only at $C=-N$ is not legitimate.
An exact formula specifying that additional data follows from the
same Mellin continuation. Put
$f(t)=\Li_A(Xe^{-t})\Li_B(Ye^{-t})$ and $f_j=[t^j]f(t)$. Then
\begin{align}\label{tc:normal-derivative}
 \left.\partial_C\mathcal T(A,B,C;X,Y)\right|_{C=-N}
 =(-1)^NN!\Bigg\{&
  \int_0^1t^{-N-1}\left(f(t)-\sum_{j=0}^Nf_jt^j\right)\dd t
  +\int_1^\infty t^{-N-1}f(t)\dd t\notag\\
 &+\sum_{j=0}^{N-1}\frac{f_j}{j-N}
        -\psi(N+1)f_N\Bigg\}.
\end{align}
Both integrals converge ordinarily. To prove this, use the Taylor
subtraction in Lemma~\ref{tr:T-entire} and
\[
 \frac1{\Gamma(-N+v)}
   =(-1)^NN!v\{1-\psi(N+1)v+O(v^2)\}.
\]
The pole $f_N/v$ supplies the last term in
\eqref{tc:normal-derivative}. The formula displays exactly what remains
after the negative-order value reduces to products of polylogarithms.
It is an integral representation for the transverse derivative, not a
claim that this derivative reduces to the value at $C=-N$.

\subsection{An ordinary three-factor log-Gamma integral}
Define the centered log-Gamma function and a constant by
\begin{equation}\label{tc:centered-Gamma}
 \Lambda(x)=\log\Gamma(x)-\tfrac12\log(2\pi)=\zeta'(0,x),
 \qquad \kappa=\gamma+\log(2\pi).
\end{equation}
Its periodic version is integrable to every finite power near an
endpoint. The identity below therefore requires no finite part.
For a variable $A$ and an integer $p\geq1$, use the operators
\begin{equation}\label{tc:Doperators}
 D_{+,A}^{(p)}=\partial_A+\log p-\kappa-\frac{i\pi}{2},
 \qquad
 D_{-,A}^{(p)}=\partial_A+\log p-\kappa+\frac{i\pi}{2}.
\end{equation}

\begin{theorem}[Unequal-frequency centered Gamma correlation]
\label{tc:Gamma-triple}
For the disjoint grids \eqref{tr:grids},
\begin{align}
 &\int_0^1\prod_{j=1}^3\Lambda(\{p_jx+a_j\})\dd x
  \notag\\
 &\quad=-\frac1{4\pi^3}\Im
   \sum_{h_1=0}^{p_1-1}\sum_{h_2=0}^{p_2-1}\sum_{h_3=0}^{p_3-1}
   \sum_{k=1}^3
   \left.
       D_{+,A}^{(p_i)}D_{+,B}^{(p_j)}D_{-,C}^{(p_k)}
       \mathcal T(A,B,C;X_k,Y_k)
   \right|_{A=B=C=1}.
   \label{tc:Gamma-formula}
\end{align}
Here $i<j$ are the other indices, and $X_k,Y_k$ are formed from the
lifted triple $(b_{1,h_1},b_{2,h_2},b_{3,h_3})$ as in
\eqref{tr:Sdef}. The imaginary part is taken only after specialization
to these real spectral orders.
\end{theorem}
\begin{proof}
Lerch's Gamma identity in \eqref{tc:centered-Gamma} gives
\[
 \int_0^1\prod_j\Lambda(\{p_jx+a_j\})\dd x
  =-\left.\partial_{\lambda_1}\partial_{\lambda_2}
               \partial_{\lambda_3}
          K_{\boldsymbol p}(\boldsymbol\lambda;\boldsymbol a)
           \right|_{\boldsymbol\lambda=(1,1,1)}.
\]
Differentiation under the integral is justified, near each separated
endpoint, by a majorant of the form
$t^{-\rho}(1+|\log t|^3)$ with a fixed $0<\rho<1$ in a sufficiently
small neighborhood of that spectral point. Insert
\eqref{tr:Kp} and \eqref{tr:unit-K}. At $\lambda=1$,
\[
 p^{\lambda-1}\Gamma(\lambda)(2\pi)^{-\lambda}=(2\pi)^{-1},
 \qquad
 \partial_\lambda\log\{p^{\lambda-1}\Gamma(\lambda)
                  (2\pi)^{-\lambda}\}=\log p-\kappa.
\]
The phase in a region with two positive frequencies is $-i$ at this
point. Its derivatives contribute $-i\pi/2$ in the two positive
positions and $+i\pi/2$ in the negative one, giving the operators
\eqref{tc:Doperators}. The opposite region is the complex conjugate
after evaluating at real orders. Thus its pair contributes
$2\Re(-iW)=2\Im W$. The common factor $(2\pi)^{-3}$ and the initial
minus sign give \eqref{tc:Gamma-formula}.
\end{proof}

The formula is a finite Gamma--polylogarithm identity with ordinary
integrals and entire spectral derivatives. The correlation of three
uncentered $\log\Gamma$ factors follows by expanding
$\log\Gamma=\Lambda+\tfrac12\log(2\pi)$. Since
$\int_0^1\Lambda(\{px+a\})\dd x=0$, the only additional terms are
the three ordinary centered bilinear correlations and
$\tfrac18\log^3(2\pi)$. Those bilinear correlations already have
single-polylogarithm order-derivative evaluations in the incoming
dilation report \cite{incoming-sh}.

\subsection{A twelfth-root digamma example with explicit lower terms}
For pairwise distinct $b_1,b_2,b_3$ modulo one, let
\[
 \mathcal P(b_1,b_2,b_3)
  =\FP_x\int_0^1\prod_{j=1}^3\psi(\{x+b_j\})\dd x.
\]
The unit-frequency formula \eqref{tr:triple-coeff} gives
\begin{equation}\label{tc:unit-digamma}
 \mathcal P(\boldsymbol b)
  =-2\Re\sum_{k=1}^3
    \left.
      D_{+,A}^{(1)}D_{+,B}^{(1)}D_{-,C}^{(1)}
        \mathcal T(A,B,C;X_k,Y_k)
    \right|_{A=B=C=0}.
\end{equation}
This is the already established translated three-digamma coordinate
formula; our example uses it after unequal-frequency transport.

Choose
\begin{equation}\label{tc:example-data}
 (p_1,p_2,p_3)=(1,2,3),\qquad
 (a_1,a_2,a_3)=(0,\tfrac13,\tfrac14).
\end{equation}
The lifted sets are
\[
 \mathcal B_1=\{0\},\quad
 \mathcal B_2=\{\tfrac16,\tfrac23\},\quad
 \mathcal B_3=\{\tfrac1{12},\tfrac5{12},\tfrac34\}.
\]
They are disjoint, and every color in \eqref{tc:unit-digamma} is a
twelfth root of unity distinct from one. Theorem~\ref{tr:finite-transport},
with $m_1=m_2=m_3=0$ and $\gamma_0=-\psi$, gives the exact identity
\begin{align}
 &\FP_x\int_0^1\psi(x)\psi(\{2x+\tfrac13\})
                              \psi(\{3x+\tfrac14\})\dd x\notag\\
 &\quad=\frac16\sum_{b\in\mathcal B_2}\sum_{c\in\mathcal B_3}
                     \mathcal P(0,b,c)+\mathcal L,\label{tc:example}\\
 \mathcal L
 &:=\frac{\log2}{3}\sum_{c\in\mathcal B_3}J_{00}(c)
       +\frac{\log3}{2}\sum_{b\in\mathcal B_2}J_{00}(b).
       \label{tc:example-lower}
\end{align}
The last constant reduces completely. With $u=\log2$, $v=\log3$,
and $\gamma_1=\gamma_1(1)$,
\begin{align}\label{tc:lower-polynomial}
 \mathcal L={}&2(\gamma_1-\zeta(2))(u+v)
   -\gamma(6u^2+4uv+3v^2)\notag\\
 &-7u^3-7u^2v-4uv^2-\tfrac32v^3.
\end{align}
To verify the reduction, Stieltjes multiplication gives
\begin{align*}
 \gamma_0(\tfrac14)+\gamma_0(\tfrac34)&=2\gamma+6\log2,\\
 \gamma_1(\tfrac14)+\gamma_1(\tfrac34)
      &=2\gamma_1-6\gamma\log2-7\log^22,\\
 \gamma_0(\tfrac13)+\gamma_0(\tfrac23)&=2\gamma+3\log3,\\
 \gamma_1(\tfrac13)+\gamma_1(\tfrac23)
      &=2\gamma_1-3\gamma\log3-\tfrac32\log^23.
\end{align*}
Apply the same multiplication law to the two lifted grids in
\eqref{tc:example-lower} and insert \eqref{tr:J00}. Collecting terms
gives \eqref{tc:lower-polynomial}.

Independent subtracted quadrature gives the following diagnostic values:
\[
 \FP_x\int_0^1\psi(x)\psi(\{2x+\tfrac13\})
                \psi(\{3x+\tfrac14\})\dd x
       \approx77.3418540125851504884052200711,
\]
\[
 \mathcal L\approx-23.0290680355012727975296691002.
\]
These are numerical illustrations of the proved identities. Their
normalization is the stated ambient finite part. In particular, removing
the lower terms in \eqref{tc:example-lower} would change the value by
the displayed nonzero constant.

% END sections/06_correlations.tex

% BEGIN sections/07_audit_verification.tex
\section{Proof status, audit findings, and reproducible verification}
\label{sec:audit}

\subsection{What has been proved, and what is inherited}
The distinction between a proved identity, a finite representation,
and a numerical candidate is essential to this continuation. The
following table records the status of its principal components.

\begin{longtable}{@{}>{\raggedright\arraybackslash}p{.26\linewidth}
>{\raggedright\arraybackslash}p{.19\linewidth}
>{\raggedright\arraybackslash}p{.49\linewidth}@{}}
\toprule
Result & Status here & Precise scope\\
\midrule
\endhead
Dougall summation and the zero-excess binomial constant &
Classical input &
Dougall supplies \eqref{dg:sum}; the constant
\eqref{dg:quarticN0} is already in the Ramanujan--B\"uhring literature.\\
\addlinespace
Centered coefficients and every integer resonance &
Proved continuation &
Theorems~\ref{dg:master}, \ref{dg:closed}, and \ref{dg:alljets}
hold in the stated positive parameter domain, including residue zeros.\\
\addlinespace
Harmonic moment identities and primitives &
Proved specializations &
All displayed sums are single convergent brackets. The primitive is
coefficientwise regular at moving lower Gamma zeros.\\
\addlinespace
Nonlinear Stieltjes coordinate change &
Proved continuation &
Theorems~\ref{nc:universal}, \ref{nc:composition},
\ref{nc:Stieltjes}, and \ref{nc:rigidity} supply all local coefficients
and their sharp vanishing range.\\
\addlinespace
Three unequal integer frequencies &
Proved continuation &
Theorem~\ref{tr:finite-transport} applies to disjoint singular grids.
Its finite reduction includes all prescribed coordinate terms.\\
\addlinespace
Weighted colored kernels and Gamma correlations &
Proved consequences &
Finite root filtering and \eqref{tc:Gamma-formula} give explicit
identities; no minimal-basis or independence claim follows.\\
\addlinespace
Current Gaussian $S_6$ and $S_8$ evaluations &
Remain conjectural &
No exact functional certificate for either short evaluation is
provided here. Their inherited numerical proximity is not equality.\\
\bottomrule
\end{longtable}

The first three source questions identified in Section~\ref{sec:scope}
are answered under explicit hypotheses. The broader collision problem,
minimal arithmetic bases for spectral derivatives, and the Gaussian
evaluation conjectures remain separate questions. The current manuscript
already treats $S_4$ as proved; it must not be reset to conjectural status
when material from older arrivals is integrated \cite{repo}.

\subsection{A verified editorial correction already proposed upstream}
The canonical integration chapter contains the wording
\emph{Basis atoms must be $\Q$-independent} in its discussion of PSLQ
\cite[chapters/07-integration.tex]{repo}. The incoming Gauss--Hurwitz
report already objects to that wording \cite[Section 8]{incoming-gh}.
That correction should be retained, with credit to the earlier report.
An appropriate replacement is:
\begin{quote}
Known rational dependencies among the supplied constants should be
removed or modeled explicitly. Otherwise PSLQ may return a pre-existing
relation with zero target coefficient instead of an evaluation of the
target.
\end{quote}
This expresses the practical requirement without imposing a prior proof
of rational independence on a discovery algorithm. A relation among
the supplied atoms is a valid output of an integer-relation procedure
\cite{pslq};
it simply may not evaluate the additional target. The change is
editorial and logical, not a new counterexample to a theorem.

\subsection{Normalization points exposed by the new proofs}
No new false theorem was found in the particular incoming results used
as premises. The following are precise implementation and integration
requirements established by this article.

\begin{enumerate}
\item \textbf{Differentiate the deformation before deleting zeros.}
Equation~\eqref{dg:revivalpi} gives an exact diagnostic:
$r_0(-1)=0$ but its first spectral coefficient is $\pi^2/8$.
The regularized Gamma product or the finite polynomial shift must be
formed before taking a resonant derivative.

\item \textbf{Specify the sum and its reference coordinate.}
The subtraction in \eqref{dg:RK} uses $n+a/2$. Replacing this reference
without the finite transport \eqref{dg:reference} changes the
counterterms. In the correlation problem, the prescribed ambient
finite part gives the nonzero scalar term in
\eqref{tr:gamma-lift}; it cannot be replaced by a
zero-mean pullback without its coordinate correction.

\item \textbf{Distinguish a pointwise derivative from a distributional
derivative.} The nonlinear coefficient for
$T_{n,r}=\FP\gamma_n^{(r)}$ is $\alpha_{n,r}$, whereas that for
$D^rG_n$ is $b_{n,r}$. Their difference is the exact contact term
\eqref{nc:derivative-contact}. The derivative comparison uses the
common smooth completion of $\gamma_n(1+x)$ across the endpoint.
Extending that entire smooth remainder by zero would introduce
additional step-function derivatives and define a different object.

\item \textbf{Keep the transformed germ in composition.}
The last term of \eqref{nc:cocycle} is
$\mathfrak A_f[h\circ g]$. Leaving $h$ unchanged there would
discard the logarithmic and principal-part coefficients generated
by the first coordinate change.

\item \textbf{Separate a negative-order value from its transverse
derivative.} Equation~\eqref{tc:weighted-negative} may be
differentiated in $A$ and $B$. A derivative in $C$ needs additional
information, supplied for the ordinary kernel by
\eqref{tc:normal-derivative}.

\item \textbf{Do not silently merge singular grids.}
The full unequal-frequency theorem assumes disjoint singular sets.
At collisions, products have higher pole order and the spectral
integral has a different convergence domain. The coordinate law for
a specified coincident product in Corollary~\ref{nc:coincident}
does not by itself identify the collision limit of separated
finite parts.
\end{enumerate}

These requirements do not identify hypothetical source errors. They
state which objects the present proofs evaluate and which tempting
changes would evaluate something else.

\subsection{Finite verification and its evidentiary role}
The supplied Python programs use exact SymPy algebra and mpmath
arithmetic. The recorded runs use SymPy~1.14.0 and mpmath~1.3.0.
Their machine-readable results are supplied with the sources.

\paragraph{Dougall and harmonic branch.}
\src{code/verify_dougall.py} checks the coefficient differential
law and residue polynomial exactly through coefficient index two,
as well as the explicit quartic $C_1$. Its 19 numerical comparisons
include generic resonances through $N=3$, first and second resonant
jets, a zero-residue case, the revived finite head, quartic moments,
squared-binomial moments and harmonic sums, and
\eqref{dg:quarteridentity}. The run uses 75 decimal digits, the first
120 summands, and a finite centered asymptotic tail through $C_{11}$.
The largest observed absolute discrepancy is
$3.76169\times10^{-36}$, below the configured $10^{-31}$ threshold.
The asymptotic-tail calculation is a diagnostic approximation; it
does not furnish a rigorous remainder enclosure.

A separate unaccelerated check of \eqref{dg:quarticN1jet} gives
partial-sum errors of about $9.41\times10^{-6}$ at 100 terms and
$2.03\times10^{-10}$ at 30,000 terms against the right side
\[
-0.0147807743355189137009913877658542374\ldots.
\]
This checks the summand directly without reusing the tail completion.
The exact proof remains Theorem~\ref{dg:alljets}.

\paragraph{Nonlinear coordinate branch.}
\src{code/verify_nonlinear.py} records 138 exact finite-algebra
checks of the coefficient generator, cocycle polynomial,
low-order formulas, and filtration. The independent implementation
\src{code/independent_cutoff.py} constructs the inverse coordinate
series, transforms the test function with its inverse Jacobian,
integrates each singular monomial in the source coordinate, and
extracts the constant after composing the primitive. Its actual-cutoff
side does not use the proposed residue formula or the
$\alpha_{n,r}$ generator.

For
\[
 f(t)=q\left(t+\frac23t^2-\frac15t^3+\frac17t^4\right),
 \qquad q=1,2,
\]
it checks 48 pairs of Stieltjes and derivative indices, through
$r=3$ and $n=2r+2$, against 140 monomial test pairings.
The $q=2$ comparisons retain exact polynomial expressions in
$\log2$. Every pairing also agrees with the reusable
coordinate-coefficient implementation. These tests independently
verify signs, Jacobians, and scale terms in low degrees.

\paragraph{Unequal-frequency branch.}
\src{code/verify_transport.py} records four exact polynomial
integrals and 13 numerical comparisons at 45 decimal digits.
The numerical checks include pointwise Stieltjes multiplication,
the nonzero ambient means, direct finite-part quadrature at
$(p_1,p_2,p_3)=(1,2,3)$, permutation symmetry, the simplified
twelfth-root lower term, and the six-region Fourier normalization.
For the latter, at orders $(2,2,2)$ the Hurwitz factors are
Bernoulli polynomials. Direct piecewise polynomial integration
for the recorded shifts and frequencies gives exactly
\[
 -\frac{1879}{125411328}.
\]
A separate Mellin evaluation of the colored kernels agrees with it.
The largest recorded scaled numerical discrepancy is about
$1.46\times10^{-45}$.

\paragraph{A numerical differentiation issue.}
During verification, generic numerical differencing of
$\zeta(s,1/4)$ at $s=0$ for the second order derivative lost
accuracy in the selected backend. Native Hurwitz spectral
derivatives, corroborated by Cauchy-circle differentiation, agreed
with the independent harmonic sum. The supplied program therefore
uses the native spectral derivative interface. This is a
computational observation from the recorded run, not a correction
to the analytic identity or a claim about every implementation.

The exact all-index results are proved by finite asymptotic
remainders, normal convergence, analytic coefficient extraction,
residues, and finite root filtering. Finite symbolic tests and
floating-point agreement provide separate implementation evidence.
No numerical file in this package is labeled an interval certificate.

\subsection{Proposed repository integration}
The modular TeX sources use the label prefixes \src{dg:},
\src{dp:}, \src{nc:}, \src{tr:}, and \src{tc:}. The source manifest
pins the inspected baseline and records the incoming archive hashes.
The companion \src{INTEGRATION_NOTES.md} gives a source-to-result
crosswalk and proposed insertion points.

A suitable integration order is: retain the incoming Gauss theorem,
append the centered Dougall extension and harmonic jets; place the
primitive beside the existing Lerch and integration material; extend
the affine finite-part section with the nonlinear law; and append the
three-frequency transport after the existing pair and translated
triple formulas. Each insertion should preserve its explicit
subtraction convention and proof status. The source reports remain
provenance, while the shared lemmas need appear only once in the
consolidated manuscript.

% END sections/07_audit_verification.tex

% BEGIN sections/08_research_questions.tex
\section{Further research questions and concrete next targets}
\label{sec:questions}

The proved results make the next questions more specific. The following
agenda concerns exact identities, structural reductions, and proof
certificates; it does not depend on searching for sharper inequalities.
Questions are labeled as questions, rather than as unproved theorems.

\subsection{Additional identities within the Dougall family}

\paragraph{1. Closed polynomial formulas for all resonant counterterms.}
The recurrence \eqref{dg:recurrence} is already finite and exact, and
\eqref{dg:quarticcoefficienttransport} builds every quartic moment
from one base sequence. Can its coefficients be written in a useful
closed combinatorial form in $N,j$, perhaps through central factorial
numbers and Bernoulli polynomials? A useful outcome would provide both
the subtraction polynomial and a compact exact certificate, avoiding
large expanded rational polynomials. This is a representation question;
the existence and the resonance value are proved here.

\paragraph{2. A minimal arithmetic presentation for rational-parameter jets.}
At rational $a,b,c$, Theorem~\ref{dg:alljets} gives a finite list of
Hurwitz derivatives and Stieltjes constants. Finite Fourier transforms
then express rational-shift Hurwitz values in residue-class or
Dirichlet-character coordinates. Which symmetries of a particular
Dougall family eliminate entire character sectors? The first concrete
case is the quartic family at shift $1/4$, comparing the Stieltjes
constants $\gamma_m(1/4)$ with the zero-order spectral jets
$\zeta^{(m)}(0,1/4)$ from its half-resonance. An exact reduction matrix
for each prescribed finite jet order would be useful. Numerical
independence is a separate question.

\paragraph{3. Mixed parameter jets along residue-zero loci.}
Formula~\eqref{dg:EN} makes all zeros explicit, but the article mainly
illustrates derivatives in $u$. When $b$ or $c$ is a positive integer
at most $N$, what short identities arise from mixed derivatives
transverse and tangent to $\rho_N=0$? The first target is a
two-parameter table at $b=c=1$, where the zero has higher multiplicity.
The finite product predicts which derivative orders first survive.
A complete statement should retain the small-index summands before
specialization, just as in \eqref{dg:revivalpi}.

\paragraph{4. Primitive endpoint values beyond the counterterms.}
The weighted primitive in Theorem~\ref{dp:primitive} lowers the
hypergeometric order and gives a finite colored-polylogarithm
description of the subtractions. It does not evaluate every endpoint
of the remaining \({}_4F_3\). Which parameter choices admit an
additional summation, a modular parametrization, or a finite
iterated-integral reduction? The squared-binomial branch is a
natural first case because its base generating function has
classical integral representations. Any proposed polylogarithmic
reduction should specify its function class and branches, rather
than infer it from the appearance of polylogarithmic counterterms.

\paragraph{5. Other exactly summable well-poised families.}
The present proof uses two special structures: reflected Gamma
shifts giving even powers, and an exact Gamma-product endpoint.
Which further very-well-poised summations retain both structures
under a one-parameter excess deformation? A precise first task is
to identify a higher-order family, derive its centered finite
asymptotic coefficients, and determine whether a differential
transport identity simplifies every resonant constant as in
\eqref{dg:closedeq}. The existence of the Dougall result does not
establish such a theorem for arbitrary hypergeometric families.

\subsection{Coordinate algebra and collisions}

\paragraph{6. Collision limits versus specified coincident finite parts.}
Let two translated singular points approach each other.
Corollary~\ref{nc:coincident} evaluates the coordinate change of the
finite part of the pointwise product at coincidence, while
Theorem~\ref{tr:finite-transport} treats disjoint supports.
What explicit local terms must be subtracted from the separated
family to obtain that coincident finite part? The first nontrivial
target is $\gamma_0(x)\gamma_0(x+\varepsilon)$, followed by
$\gamma_m^{(r)}(x)\gamma_n^{(s)}(x+\varepsilon)$.
The regulator and the relative approach direction must be fixed
before taking the limit. Equality of two unspecified
regularizations would not be a meaningful claim.

\paragraph{7. An intrinsic classification of the coordinate corrections.}
The residue law \eqref{nc:cocycle} is explicit. Can all local
corrections compatible with scalar pullback, finite principal parts,
and composition be classified, with the freedom of adding fixed
delta derivatives made precise? The sharp filtration
\eqref{nc:order-filtration} provides a concrete family of graded
invariants for that classification. A first step is to classify
the finite jet action through pole order three and determine
which coefficients can be changed by a fixed normalization.

\paragraph{8. Smooth coordinates and noninteger logarithmic exponents.}
For fixed pole order, the proof only uses finite coordinate jets
and an integrable remainder. This suggests an extension from
analytic coordinates to a stated finite differentiability class.
The exact minimum regularity and remainder hypotheses should be
proved, not inferred from formal Taylor manipulations. A distinct
question is the extension of $\log^n x$ to noninteger powers of
$\log x$, where branch choices and the polynomial filtration change.
It should be treated separately from the integer Stieltjes hierarchy.

\subsection{Colored kernels and exact relation certificates}

\paragraph{9. Minimal finite reductions for three or more frequencies.}
The finite lift has at most $p_1p_2p_3$ translated triples, and each
triple has six sign regions. Root filtering establishes existence
of a finite ordinary-kernel presentation. Can common divisors,
translation symmetry, and color conjugation compress this
presentation canonically? The first target is an exact reduction
algorithm for the $(1,2,3)$ example, with the order-derivative
coordinates retained. For four frequencies, one should specify
the resulting triple-sum kernels and a finite sign-region
decomposition before claiming analogous closure.

\paragraph{10. Transverse derivatives at negative Tornheim order.}
At $C=-N$, the kernel reduces to products of ordinary
polylogarithms. Equation~\eqref{tc:normal-derivative} isolates the
extra information in the first normal derivative. At roots of
unity, which choices of $A,B,N$ reduce that integral to a finite
combination of known depth-two polylogarithms, order derivatives,
and Gamma or Stieltjes data? The case $C=0$ with $A,B$ equal to
zero or one is a concrete starting point. The ordinary
negative-order value alone contains insufficient information.

\paragraph{11. Exact Gaussian $S_6$ and $S_8$ certificates.}
The retained conjectures need an exact relation in a specified
functional or iterated-integral system, together with endpoint
control and evaluation. The appropriate first task is to pin one
candidate vector, its alphabet, and the exact relation generators.
Either produce a certificate whose evaluation is the desired
identity, or prove exclusion from that \emph{specified} generated
relation space. Exclusion from a selected formal space would
identify a missing mechanism; it would not prove numerical
independence or disprove the period identity. The current source
ledger already emphasizes this distinction \cite{repo}.

\paragraph{12. A common exact coefficient and certificate format.}
The Bernoulli recurrence, the nonlinear residue, and the finite root
filter are all constructive operations on finite data once the
analytic theorem has been fixed. A useful ProveIt implementation
would record the theorem identifier, parameter domain, local
coordinate, coefficient order, finite-head terms, and exact
polynomial certificate in one reusable format. A first deliverable
could replay \eqref{dg:quarticN1jet},
\eqref{nc:r2unit}, and \eqref{tc:lower-polynomial} using
independent exact arithmetic. Such a format would support
integration and auditing without treating a floating-point
comparison as a proof certificate.

\subsection{Research priorities}
The shortest paths to additional exact identities appear to be
mixed Dougall jets at explicit residue zeros, arithmetic reduction
of the rational-shift special values, and normal derivatives of
colored Tornheim kernels at $C=0$. Collision limits require a
separate normalization theorem. The Gaussian conjectures remain
the most direct targets for a finite exact relation certificate,
but the present continuation supplies no reason to mark them
proved before that certificate exists.

% END sections/08_research_questions.tex

% BEGIN references.tex
\addcontentsline{toc}{section}{References}
\begin{thebibliography}{99}
\interlinepenalty=10000

\bibitem{repo}
V. Reshetnikov, \emph{ProveIt}, consolidated polylogarithm manuscript,
editorial ledger, and incoming research archives. Snapshot
\pin, inspected 10 October 2026.
\url{https://github.com/VladimirReshetnikov/ProveIt/tree/445f754610e1377939235794de84ed9075fc09f5/Analysis/Polylogarithms/docs/manuscript}.
Incoming directory:
\url{https://github.com/VladimirReshetnikov/ProveIt/tree/445f754610e1377939235794de84ed9075fc09f5/docs/incoming}.
The companion source manifest records the exact files and Git blob hashes.

\bibitem{incoming-gh}
ProveIt incoming research report,
\emph{Gauss--Hurwitz Pole Cancellation: Resonant gamma-ratio sums,
harmonic identities, Stieltjes jets, and polylogarithmic endpoint
subtraction}, 10 October 2026.
Archive \src{ProveIt_Gauss_Hurwitz_2026-10-10.zip} in \cite{repo}.
Especially Sections 2--7, the audit in Section 8, and the Dougall
question in Section 9.

\bibitem{incoming-md}
ProveIt incoming research report,
\emph{Mellin--Lerch Transforms and Dilation Identities for Stieltjes
Correlations: Polylogarithms, generalized harmonic sums, and exact
finite parts}, 10 October 2026.
Archive \src{ProveIt_Mellin_Dilation_2026-10-10.zip} in \cite{repo}.
Especially the dilation section and the nonlinear-coordinate question
in \src{sections/06-audit-research.tex}.

\bibitem{incoming-sh}
ProveIt incoming research report,
\emph{Exact Stieltjes and Harmonic Identities: Integer dilations,
three-factor products, Gamma coefficients, and corrected conjectures},
10 October 2026.
Archive \src{ProveIt_Stieltjes_Harmonic_Identities_2026-10-10.zip}
in \cite{repo}. Especially \src{sections/dilation.tex},
\src{sections/trilinear.tex}, and \src{sections/further_research.tex}.

\bibitem{incoming-nhj}
ProveIt incoming research report,
\emph{Nested Harmonic Jets and Gamma Renormalization},
10 October 2026.
Archive \src{ProveIt_Nested_Harmonic_Jets_2026-10-10.zip} in \cite{repo}.
Inspected for the harmonic generating calculus and inherited proof status.

\bibitem{incoming-continuation}
ProveIt incoming research report,
\emph{Further Identities and Critical Transitions in the
Polylogarithm--Stieltjes Calculus: Translation geometry, Euler reversal,
Lerch velocities, and depth in the Cayley shuffle quotient},
10 October 2026.
Archive \src{ProveIt_Polylogarithm_Stieltjes_Continuation_2026-10-10.zip}
in \cite{repo}. The scope and relevant audit and questions were checked.

\bibitem{dlmf-hyper}
NIST Digital Library of Mathematical Functions,
\emph{Section 16.4: Argument Unity}, especially the Rogers--Dougall
very-well-poised summation, Eq.~16.4.9.
\url{https://dlmf.nist.gov/16.4#E9}.

\bibitem{dlmf-gamma}
NIST Digital Library of Mathematical Functions,
\emph{Section 5.11: Asymptotic Expansions}, especially Eq.~5.11.8
and the Gamma-ratio expansions in Section 5.11(iii).
\url{https://dlmf.nist.gov/5.11}.

\bibitem{dlmf-hurwitz}
NIST Digital Library of Mathematical Functions,
\emph{Section 25.11: Hurwitz Zeta Function}.
Fourier continuation, functional relations, multiplication,
Bernoulli values, and the Gamma identity, Eq.~25.11.18.
\url{https://dlmf.nist.gov/25.11}.

\bibitem{dlmf-lerch}
NIST Digital Library of Mathematical Functions,
\emph{Section 25.14: Lerch's Transcendent}.
\url{https://dlmf.nist.gov/25.14}.

\bibitem{buhring}
W. B\"uhring,
\emph{Partial sums of hypergeometric series of unit argument},
Proceedings of the American Mathematical Society \textbf{132}
(2004), 407--415.
\url{https://arxiv.org/abs/math/0311126}.
Section 4, Corollaries 1--2, records the classical zero-excess
constants and their Ramanujan--Berndt--Evans antecedents.

\bibitem{em-hurwitz}
O. Espinosa and V. H. Moll,
\emph{On some integrals involving the Hurwitz zeta function: Part 1},
The Ramanujan Journal \textbf{6} (2002), 159--188.
\url{https://www.math.tulane.edu/~vhm/papers_html/hurwitz1.pdf}.

\bibitem{em-tornheim}
O. Espinosa and V. H. Moll,
\emph{The evaluation of Tornheim double sums. Part 1},
Journal of Number Theory \textbf{116} (2006), 200--229.
\url{https://arxiv.org/abs/math/0505647}.

\bibitem{zhao}
J. Zhao,
\emph{A note on colored Tornheim's double series},
Integers \textbf{10} (2010), A59, 879--882.
\url{https://arxiv.org/abs/0907.5106}.

\bibitem{pslq}
H. R. P. Ferguson and D. H. Bailey,
\emph{A polynomial time, numerically stable integer relation algorithm},
RNR Technical Report RNR-91-032, revised 14 July 1992.
\url{https://www.davidhbailey.com/dhbpapers/pslq.pdf}.

\end{thebibliography}

\end{document}
